% !TEX program = xelatex
% 马克思主义基本原理 · 考点边界清理与完整讲解
% 依据：高等教育出版社《马克思主义基本原理》（马克思主义理论研究和建设工程重点教材，2023 年版）章节体系
% 2026-11-16 生成；源资料：学习者提供的 15 单选 + 5 多选原始作答、留空与订正记录
\documentclass[11pt]{ctexart}
%=====================================================================
% theme-violet.tex —— 讲义精讲 · 紫花主题
% 用法：\documentclass[11pt]{ctexart}
%       \input{assets/theme-violet}
%       \begin{document} ... \end{document}
% 只用 ctexart 自带能力 + xcolor/enumitem/fancyhdr/pifont/amsmath，
% 不依赖 tikz/tcolorbox，保证 MiKTeX 可直接编译。
%=====================================================================

\usepackage[a4paper,left=2.0cm,right=2.0cm,top=2.1cm,bottom=2.3cm,%
            headsep=12pt,footskip=22pt]{geometry}
\usepackage{amsmath,amssymb}
\usepackage{needspace}
\usepackage{xcolor}
\usepackage{array}
\usepackage{enumitem}
\usepackage{pifont}
\usepackage{fancyhdr}
\usepackage[hidelinks]{hyperref}

%---------------------------------------------------------------- 配色
\definecolor{violet}{RGB}{107,78,143}      % 主色：紫
\definecolor{violetdeep}{RGB}{63,44,92}    % 深紫：标题/强调
\definecolor{petal}{RGB}{183,159,208}      % 花瓣：装饰线、花饰
\definecolor{bandviolet}{RGB}{247,243,251} % 浅紫底：档案/提示框
\definecolor{bandcream}{RGB}{252,249,242}  % 米白底：易错/注意框
\definecolor{rulepurple}{RGB}{217,207,230} % 分隔线
\definecolor{inkgray}{RGB}{92,96,110}      % 次要文字

%---------------------------------------------------------------- 字体
% 有 Noto CJK 用 Noto，没有就退回 ctex 默认（Windows 下为宋体/黑体）
\IfFontExistsTF{Noto Serif CJK SC}{\setCJKmainfont{Noto Serif CJK SC}}{}
\IfFontExistsTF{Noto Sans CJK SC}{\setCJKsansfont{Noto Sans CJK SC}}{}

%---------------------------------------------------------------- 版式
\setlist{nosep,leftmargin=1.9em}
\linespread{1.28}
\setlength{\parskip}{3.5pt}
\setlength{\parindent}{0pt}
\setlength{\tabcolsep}{5pt}
\renewcommand{\arraystretch}{1.28}

%---------------------------------------------------------------- 花饰
% \ding{95} = 五瓣小花，\ding{96} = 星形花；都来自 pifont，必定存在
\newcommand{\orn}{\textcolor{petal}{\ding{95}}}          % 行内小标
\newcommand{\ornstar}{\textcolor{petal}{\ding{96}}}      % 星形花
\newcommand{\flowerrule}{%                                % 居中花饰分隔线
  \par\vspace{7pt}%
  \noindent\makebox[\linewidth]{\color{petal}%
    \rule{0.20\linewidth}{0.6pt}\hspace{8pt}{\small\ding{95}}\hspace{4pt}%
    \raisebox{1.1pt}{\tiny\ding{95}}\hspace{4pt}{\small\ding{95}}\hspace{8pt}%
    \rule{0.20\linewidth}{0.6pt}}%
  \par\vspace{7pt}}

%---------------------------------------------------------------- 标题与节头
\newcommand{\doctitle}[3]{%   #1 主标题  #2 副标题  #3 元信息（可用 \\ 分行）
  \begin{center}
    {\color{petal}\small \orn\enspace #2\enspace \ornstar}\\[7pt]
    {\LARGE\bfseries\color{violetdeep} #1}\\[10pt]
    {\color{petal}\rule{0.55\linewidth}{1pt}}\\[4pt]
    {\small\color{inkgray} #3}
  \end{center}\vspace{2pt}}

\newcommand{\hdr}[2]{%   #1 章节号（如 §0）  #2 标题
  \par\vspace{12pt}%
  \noindent\colorbox{violetdeep}{\color{white}\bfseries\small\;#1\;}\hspace{7pt}%
  {\color{violetdeep}\large\bfseries #2}\par\nobreak
  \vspace{3pt}{\color{petal}\hrule height 1.2pt}\vspace{6pt}}

\newcommand{\subhdr}[1]{%
  \par\vspace{6pt}%
  \noindent{\color{violet}\bfseries #1}\par\vspace{1pt}}

%---------------------------------------------------------------- 提示框
\newcommand{\callout}[1]{% 浅紫底：讲义档案、读法说明
  \par\medskip\noindent
  \fcolorbox{rulepurple}{bandviolet}{\parbox{\dimexpr\linewidth-2\fboxsep-2\fboxrule}{\vspace{2pt}#1\vspace{2pt}}}%
  \par\medskip}

\newcommand{\warnbox}[1]{% 米白底：易错、边界条件
  \par\medskip\noindent
  \fcolorbox{rulepurple}{bandcream}{\parbox{\dimexpr\linewidth-2\fboxsep-2\fboxrule}{\vspace{2pt}#1\vspace{2pt}}}%
  \par\medskip}

\newcommand{\gapitem}[2]{% #1 缺口类型（紫） #2 内容
  \par\vspace{2pt}\noindent
  \hangindent=1.4em\hangafter=1
  {\color{petal}\small\ding{96}}\enspace{\bfseries\color{violetdeep}#1}\enspace #2\par}

% 题干块（\qstem）：每道例题／作业的「详解」开头必须先复述题干与出处。
% 起因：2026-10-25 第五章讲义把七道作业的题面集中放在一页、详解另起一页，
% 读者翻到「第 4 题」时看不到条件，等于一道没有题干的题目。
% 规则：详解块内自带题干，读者不翻页就能读懂题在问什么。
\newcommand{\qstem}[1]{%
  \par\medskip\noindent
  \fcolorbox{violet}{bandviolet}{%
    \parbox{\dimexpr\linewidth-2\fboxsep-2\fboxrule}{%
      \vspace{2pt}%
      {\small\bfseries\color{violet}题干}\hspace{0.8em}#1%
      \vspace{2pt}}}%
  \par\medskip}

% v2: 保护题干与开头，长题可跨页，续页补上下文。
\newcommand{\qbreak}{\needspace{6\baselineskip}}
% 独立练习：ID、来源、题干、考点、答案标签（通常 ans:ID）。
\newcommand{\practice}[5]{%
  \qbreak\subhdr{练习 #1}\qstem{（#2）#3}%
  \textbf{考点}：#4\par
  详解见附录第~\pageref{#5}~页。\par\medskip}
% 附录开头：ID、来源、条件摘要；自动产生 ans:ID 标签。
\newcommand{\answerstart}[3]{%
  \qbreak\subhdr{#1\quad 详解}\label{ans:#1}\qstem{（#2）#3}}
\newcommand{\answercontinue}[2]{%
  \needspace{5\baselineskip}\subhdr{#1\quad 续解}%
  {\small\color{inkgray}本页接续：#2}\par\medskip}
% 延迟提供默认值，让旧讲义可在导言区保留自己的 \newcommand{\step}。
\AtBeginDocument{\providecommand{\step}[1]{\needspace{3\baselineskip}\par\medskip\textbf{#1}\par}}

% 前提旁注：只放在「需要理解的那一步」后面，给陈述 + 本题代入，不重证。
\newcommand{\aside}[1]{%
  \par\vspace{1pt}\noindent
  \hangindent=1.6em\hangafter=1
  {\small\color{inkgray}%
    \hspace*{0.6em}{\color{petal}$\triangleright$}\hspace{2pt}#1}%
  \par\vspace{2pt}}

%---------------------------------------------------------------- 表格
% 本机 MiKTeX 的 colortbl 与 array 版本不匹配（\rowcolor 会报错），
% 故表格统一用「上粗紫线 + 表头 + 中细线」的三线式，不依赖 colortbl。
\newcommand{\tabtop}{\par\vspace{3pt}{\color{violet}\hrule height 1.2pt}\vspace{1pt}}
\newcommand{\tabmid}{\par\vspace{1pt}{\color{petal}\hrule height 0.5pt}\vspace{1pt}}
\newcommand{\tabbot}{\par\vspace{3pt}{\color{rulepurple}\hrule height 0.5pt}\vspace{8pt}}
\newcommand{\thead}[1]{{\bfseries\color{violetdeep}#1}}

%---------------------------------------------------------------- 页眉页脚
\pagestyle{fancy}
\fancyhf{}
\renewcommand{\headrulewidth}{0.4pt}
\renewcommand{\headrule}{\hbox to\headwidth{\color{rulepurple}\leaders\hrule height \headrulewidth\hfill}}
\fancyhead[L]{\small\color{inkgray}\songti \doctag}
\fancyhead[R]{\small\color{inkgray}\songti \coursetag}
\fancyfoot[C]{\small\color{inkgray}\orn\enspace{\color{violet}\thepage}\enspace\orn}

% 每篇讲义在导言区设这两个值
\newcommand{\doctag}{}
\newcommand{\coursetag}{}
\newcommand{\settags}[2]{\renewcommand{\doctag}{#1}\renewcommand{\coursetag}{#2}}

%---------------------------------------------------------------- 数学记号
\newcommand{\vect}[1]{\boldsymbol{#1}}
\newcommand{\dd}{\,\mathrm{d}}
\newcommand{\ddx}{\frac{\mathrm{d}}{\mathrm{d}x}}
\addtolength{\jot}{2pt}



\usepackage{tikz}
\usetikzlibrary{arrows.meta,calc,positioning}
\setlength{\headheight}{22pt}
\linespread{1.24}
\raggedbottom
\settags{专题 · 考点边界清理}{马克思主义基本原理}

%--------------------------------------------------------------- 记号
\newcommand{\dw}[1]{\textbf{\color{violetdeep}#1}}   % 定位词、关键判据
\newcommand{\figcap}[1]{\par\vspace{2pt}{\small\color{inkgray}#1}\par}
% 表格横线：在 tabular 内部用 \noalign，避免 \\ 后的 \trulemid 落进首列
\newcommand{\truletop}{\hline}
\newcommand{\trulemid}{\hline}
\newcommand{\trulebot}{\hline}
\renewcommand{\arraystretch}{1.35}
\setlength{\arrayrulewidth}{0.6pt}
\let\orighdr\hdr
\renewcommand{\hdr}[2]{\par\needspace{12\baselineskip}\orighdr{#1}{#2}}
\let\origsubhdr\subhdr
\renewcommand{\subhdr}[1]{\par\needspace{4\baselineskip}\origsubhdr{#1}}
\setlength{\abovedisplayskip}{5pt plus 2pt minus 2pt}
\setlength{\belowdisplayskip}{5pt plus 2pt minus 2pt}
\setlength{\emergencystretch}{1.5em}
\widowpenalty=10000
\clubpenalty=10000
\tikzset{>=Stealth}

\begin{document}

\doctitle{马原考点边界清理}
{同一批概念，为什么一到定语与定位词就掉分}
{专题讲义\quad 2026-11-16\quad 教材依据：《马克思主义基本原理》\\
高等教育出版社·马克思主义理论研究和建设工程重点教材（2023 年版）章节体系}

\hdr{§0}{导览：这本讲义为什么按「边界」组织}

你这套卷子的失分很有规律：社会存在与社会意识、生产方式、经济基础与上层建筑、资本主义再生产的实质，这些\dw{成体系的大块}你一次做对；一旦题目把两个相近概念的定语一换、把“根本观点”从辩证法搬到认识论、把三要素砍成一条，你就跟着错了。这说明问题不在“没背过”，而在\dw{背过的条目与它所在的位置对不上}——概念是散的，每一条都挂在不同的钉子上，钉子一挪，条目就乱。

所以这本讲义不按教材页码顺序复述，而是按\dw{四条边界}重排材料：它属于哪个层次（哲学／政治经济学／科学社会主义）、它的定位词是什么（源泉动力／形式状态／方向道路）、它的定语有没有被换掉（能动反映／直观反映，性质未变／性质改变）、它的固定条目数够不够（三要素、五特征、四表现）。每一节都按同一个结构走：教材怎么说 $\rightarrow$ 定位词 $\rightarrow$ 怎样判别 $\rightarrow$ 错项长什么样 $\rightarrow$ 一个小例。

\begin{center}\small
\begin{tabular}{@{}p{5.0cm}p{3.4cm}p{2.1cm}p{4.6cm}@{}}
\truletop
\thead{你原来的失分点} & \thead{教材位置} & \thead{本讲义} & \thead{处理方式}\\
\trulemid
三大规律功能混淆（原卷第 4 题） & 第一章第三节 & §3.2--3.5 & 三条定位词逐条钉死，再给对照表\\
\trulemid
矛盾观点／实践观点错位（第 6 题） & 第一章第三节、第二章第一节 & §4.1 & 两个“根本观点”分层，各归其位\\
\trulemid
“直观反映”误选（多选第 3 题） & 第二章第一节 & §4.3--4.5 & 三条认识路线横向对比\\
\trulemid
劳动力价值三要素漏项（多选第 4 题） & 第四章第二节 & §6.3 & 完整条目 + 记忆锚 + 为何是这三项\\
\trulemid
国家垄断资本主义性质误判（第 12 题） & 第五章第二节 & §7.4 & 三条论证，逐条对应“性质未变”\\
\trulemid
列宁五特征零散（多选第 5 题） & 第五章第一节 & §7.3 & 原话 + 顺序记忆锚\\
\trulemid
理论来源／组成部分混淆（多选第 1 题） & 导论 & §2.2--2.3 & 两张表交叉核对\\
\trulemid
无产阶级革命根本问题（第 13 题） & 第三章第二节（第六章同结论） & §5.4 & 命题层级辨析 + 为什么不是领导权\\
\trulemid
“两个绝不会”出处（第 15 题） & 第六章、第七章 & §8.6 & 三篇文献时间线 + 原话\\
\trulemid
货币转化为资本的前提（第 10 题留空） & 第四章第二节 & §6.2 & 因果链一步不缺\\
\trulebot
\end{tabular}
\end{center}

\subhdr{全书主线概览}

下面这张图只画路线与依赖，不代替后文的定义、原话和例子。读法：上面两个框是全书要解决的问题和四条判据；中间三个板块各自回答一个问题——哲学回答“世界怎样存在、人怎样认识它”，政治经济学回答“剩余价值从哪里来、又怎样被占有”，科学社会主义回答“这条路通向哪里、凭什么说必然”；底部是考场上真正要做的四步动作。

\begin{center}
\begin{tikzpicture}[x=1cm,y=1cm,
  ovcore/.style={draw=violetdeep,fill=bandviolet,rounded corners=3pt,align=left,inner xsep=6pt,inner ysep=4pt,line width=1pt,font=\footnotesize},
  ovbox/.style={draw=violet,fill=bandviolet,rounded corners=2.5pt,align=left,inner xsep=5pt,inner ysep=3pt,line width=0.7pt,font=\footnotesize},
  ovcream/.style={draw=violet,fill=bandcream,rounded corners=2.5pt,align=left,inner xsep=5pt,inner ysep=3pt,line width=0.7pt,font=\footnotesize},
  ova/.style={-{Stealth[length=2.2mm,width=1.6mm]},violet,line width=0.9pt},
  ovline/.style={violet,line width=0.9pt}]

\node[ovcore,text width=15.2cm,minimum height=1.62cm,anchor=north west] (CORE) at (0.25,15.35) {%
\textbf{\color{violetdeep}全书要解决的问题}\quad 概念背过了，为什么换一个定语就判错？\\
查法只有一条：把这个考点放回教材的\textbf{层次}，读出它的\textbf{定位词}，数清它的\textbf{条目}；\\
四个板块的每一个失分点，都能落到下面四项判据中的一项。};

\node[ovbox,text width=15.2cm,minimum height=2.28cm,anchor=north west] (JUDGE) at (0.25,13.42) {%
\textbf{\color{violetdeep}1\quad 定语}\quad{\color{inkgray}“能动／直观”“未变／改变”：换一个限定词，结论就换人}\\[1pt]
\textbf{\color{violetdeep}2\quad 定位词}\quad{\color{inkgray}源泉动力／形式状态／方向道路：三条各管一件事}\\[1pt]
\textbf{\color{violetdeep}3\quad 层次}\quad{\color{inkgray}辩证法／认识论／唯物史观：各管各的“根本观点”}\\[1pt]
\textbf{\color{violetdeep}4\quad 条目}\quad{\color{inkgray}三要素、五特征、四表现：多选题漏一条就零分}};

\draw[ovline] (8.1,11.14) -- (8.1,10.72);
\draw[ovline] (2.55,10.72) -- (14.05,10.72);
\draw[ova] (2.55,10.72) -- (2.55,10.15);
\draw[ova] (8.1,10.72) -- (8.1,10.15);
\draw[ova] (14.05,10.72) -- (14.05,10.15);


\node[ovbox,text width=5.0cm,minimum height=2.55cm,anchor=north west] (P1) at (0.25,10.02) {%
\textbf{\color{violetdeep}哲学·世界与人}\ {\color{inkgray}（§3--§5）}\\[1pt]
第一章\ 三大规律·五范畴\\
第二章\ 实践与能动反映\\
第三章\ 社会存在与政权\\
{\color{inkgray}关键词：源泉动力·能动·政权}};

\node[ovbox,text width=5.0cm,minimum height=2.55cm,anchor=north west] (P2) at (5.68,10.02) {%
\textbf{\color{violetdeep}政治经济学·剥削}\ {\color{inkgray}（§6--§7）}\\[1pt]
第四章\ 商品·劳动力·剩余价值\\
第五章\ 垄断·五特征·国家垄断\\
{\color{inkgray}关键词：三要素、m/v、性质未变}};

\node[ovbox,text width=5.0cm,minimum height=2.55cm,anchor=north west] (P3) at (11.11,10.02) {%
\textbf{\color{violetdeep}科学社会主义}\ {\color{inkgray}（§2、§8）}\\[1pt]
导论\ 三大来源·三大组成\\
第六章\ 空想三阶段·两大基石\\
第七章\ 两个必然·两个绝不会\\
{\color{inkgray}关键词：来源非组成·文献}};

\node[ovcream,text width=15.2cm,anchor=north west] (BOT) at (0.25,6.95) {%
\textbf{\color{violetdeep}考场四步（§1 详述，§9 有卡片）}\quad
\textbf{圈主语}（问的是规律、观点还是制度）$\rightarrow$
\textbf{找限定}（“首要”“根本”“直接”“已经”）$\rightarrow$
\textbf{数条目}（多选先默条目再选项）$\rightarrow$
\textbf{回原句}（教材怎么定位这个词）。\\[2pt]
\textbf{\color{violetdeep}文献锚点}\quad
1848《共产党宣言》“两个必然”｜1859《〈政治经济学批判〉序言》“两个绝不会”｜
1867《资本论》第一卷（剩余价值）｜1875《哥达纲领批判》（两阶段）｜
1916《帝国主义是资本主义的最高阶段》（五特征）。};
\end{tikzpicture}
\figcap{图 0-1\quad 全书主线：四条判据逐条套在三个板块上；底部是考场动作与文献锚点。}
\end{center}

\callout{\textbf{读法}\quad 第一次读不必认识图里所有词。只要记住三件事：哲学板块的三个失分点都是“定位词与层次”问题；政治经济学板块的两个失分点都是“条目与性质”问题；科学社会主义板块的失分点在“归属与出处”。带着这三件事往下读，每节的错项都能归到其中一类。}

\hdr{§1}{一条总纲：马原选择题的四种「边界」}

教材的每一个判断都是“主语 + 定语 + 谓词”的完整结构，出题人不需要编新知识，只要挪动其中一个部件，就能造出一个看起来对、实际上错的选项。你这套卷子的八处失分，按构造方式只有四类。

\begin{center}\small
\begin{tabular}{@{}p{2.7cm}p{5.2cm}p{4.4cm}p{3.5cm}@{}}
\truletop
\thead{设陷方式} & \thead{选项长什么样} & \thead{自查动作} & \thead{本次实例}\\
\trulemid
\textbf{定位词错配} & 把 A 规律的功能安到 B 规律头上 & 三大规律各背一句定位词，见“揭示”就问哪一条 & 第 4 题：把“源泉和动力”给了否定之否定\\
\trulemid
\textbf{层次错配} & 把第一章的“根本观点”搬到第二章当“首要观点” & 先问这句话在讲辩证法还是认识论 & 第 6 题：矛盾观点与实践观点之争\\
\trulemid
\textbf{定语替换} & 把“能动”换成“直观”，把“未变”换成“改变” & 圈出形容词，回教材原句核对 & 多选第 3 题 D 项；第 12 题 A 项\\
\trulemid
\textbf{条目残缺} & 多选题只给一半正确条目，另一半换成错的 & 先默写完整条目，再逐项比对 & 多选第 4、5 题\\
\trulebot
\end{tabular}
\end{center}

\subhdr{考场上真正要做的四步}

\begin{enumerate}[label=\textbf{\arabic*.},leftmargin=2.2em]
\item \dw{圈主语}：这道题问的是“规律”“观点”“制度”“前提”还是“标志”？主语决定你该去翻哪一节，而不是凭语感选。
\item \dw{找限定}：把“首要的”“基本的”“根本的”“直接的”“已经”“完全”这些词圈出来。马原的判断题，一半的分在这些词上：把“基本”读成“根本”、把“重要”读成“唯一”，就会选错。
\item \dw{数条目}：多选题先在心里默写条目（三要素？五特征？），再逐项看选项有没有偷换其中一条。漏选和错选同分，所以“少选保平安”是幻觉。
\item \dw{回原句}：拿不准时回到教材的那一句，不看辅导书的转述。教材的表述往往自带限定，例如“没有改变垄断资本主义的性质”“是资本主义生产关系的局部调整”。
\end{enumerate}

\warnbox{\textbf{一句话}\quad 马原选择题不考你“知不知道这个词”，考你“知不知道这个词挂在哪根钉子上”。钉子只有四类：层次、定位词、定语、条目数。}

从这个角度看，你原来的错法其实是同一种错法的四次重复：\dw{概念的钉子没钉牢，于是只能凭“感觉像”来选}。后面每一节的任务，就是把相应的钉子重新钉一遍，并且钉在你自己的错项上。

% ======================= PART2 =======================
\providecommand{\quoteorig}[1]{\par\smallskip\begingroup
  \setlength{\leftskip}{1.5em}\setlength{\rightskip}{1.5em plus 1em}%
  \color{violetdeep}\noindent #1\par\endgroup\smallskip}

\hdr{§2}{导论：三大来源与三大组成部分，两套“三”各回答什么}

这是你多选第 1 题涂改的地方。两套“三”都要背，但它们回答的不是同一个问题：\dw{三大来源}回答“马克思主义从哪些思想材料里长出来”，是思想史；\dw{三大组成部分}回答“马克思主义这个体系由哪几块构成”，是体系结构。把来源填进“组成部分”的空格，或者把组成部分当成“理论来源”，就是这次涂改的根源。

\subhdr{2.1\quad 创立与问世：先记住四个时间点}

马克思主义是马克思、恩格斯在 19 世纪 40 年代创立的，它吸收前人成果又实现了根本变革。四个时间点足够支撑导论部分的判断题：

\begin{itemize}[leftmargin=1.9em]
\item 1848 年 2 月，《共产党宣言》发表，\dw{标志着马克思主义公开问世}（不是“最终形成”，也不是“诞生于书斋”）。
\item 1867 年，《资本论》第一卷出版，剩余价值学说得到系统阐述。
\item 1871 年，巴黎公社革命，是无产阶级夺取政权的第一次伟大尝试（不是“成功建立社会主义”）。
\item 1883 年马克思逝世后，恩格斯整理出版《资本论》第二、三卷，并写下《反杜林论》《家庭、私有制和国家的起源》等。
\end{itemize}

\subhdr{2.2\quad 三大组成部分：体系的“三块”}

\begin{center}\small
\begin{tabular}{@{}p{3.0cm}p{4.8cm}p{3.6cm}p{4.2cm}@{}}
\truletop
\thead{组成部分} & \thead{研究什么} & \thead{在体系中的地位} & \thead{代表成果}\\
\trulemid
马克思主义哲学（辩证唯物主义与历史唯物主义） & 世界的本质、发展的一般规律，以及人认识世界和改造世界的一般方法 & 整个体系的\textbf{理论基础}（世界观和方法论） & 《关于费尔巴哈的提纲》《德意志意识形态》《反杜林论》\\
\trulemid
马克思主义政治经济学 & 资本主义生产方式及其内在矛盾、剩余价值的来源与分配 & 体系的\textbf{主要内容}，是对资本主义的解剖 & 《资本论》《剩余价值理论》\\
\trulemid
科学社会主义 & 无产阶级解放的条件、进程和目标，社会主义代替资本主义的规律 & 体系的\textbf{核心与结论}（前两者得出的结论） & 《共产党宣言》《哥达纲领批判》《社会主义从空想到科学的发展》\\
\trulebot
\end{tabular}
\end{center}

\subhdr{2.3\quad 三大理论来源：思想材料的“三支”}

\begin{center}\small
\begin{tabular}{@{}p{3.4cm}p{3.4cm}p{5.0cm}p{3.8cm}@{}}
\truletop
\thead{理论来源} & \thead{代表人物} & \thead{被批判继承的合理成分} & \thead{主要成为哪一块的材料}\\
\trulemid
德国古典哲学 & 黑格尔、费尔巴哈 & 黑格尔辩证法的“合理内核”（发展、矛盾、否定之否定），费尔巴哈唯物主义的“基本内核”（物质第一性） & 马克思主义哲学\\
\trulemid
英国古典政治经济学 & 亚当·斯密、大卫·李嘉图 & 劳动价值论：价值由劳动创造，劳动是财富的源泉 & 马克思主义政治经济学\\
\trulemid
英法空想社会主义 & 圣西门、傅立叶、欧文 & 对资本主义的批判、对未来社会的设想（消灭私有制、人人劳动、按劳分配） & 科学社会主义\\
\trulebot
\end{tabular}
\end{center}

\subhdr{2.4\quad 为什么“来源”不能被当成“组成部分”}

教材在导论里讲“马克思主义是对人类文明成果的继承与创新”，列宁则在《马克思主义的三个来源和三个组成部分》中把这两套“三”并列提出。\dw{“来源”是可以错的思想材料，“组成部分”必须是科学的理论}——这正是两栏不能互换的理由：

\begin{itemize}[leftmargin=1.9em]
\item 空想社会主义对资本主义的批判是深刻的，但它的历史观是唯心主义的（把希望寄托在天才人物和统治者的善意上），所以它是\dw{来源}，不是\dw{组成部分}；科学社会主义才是组成部分，因为它建立在唯物史观和剩余价值学说两块基石上。
\item 德国古典哲学中的黑格尔辩证法是唯心主义的，费尔巴哈唯物主义是半截子的（历史观仍是唯心主义），所以它们只能作为\dw{来源}被批判继承；马克思主义哲学是经过改造后的\dw{组成部分}。
\item 反过来，如果把“马克思主义哲学”说成“理论来源”，等于说它从自己那里长出来，逻辑上就错了。
\end{itemize}

\subhdr{2.5\quad 两大发现：社会主义为什么能从空想变为科学}

恩格斯把\dw{唯物史观}和\dw{剩余价值学说}称为马克思一生的两个伟大发现。前者揭示了社会发展的客观规律（生产力决定生产关系、经济基础决定上层建筑），后者揭示了资本主义剥削的秘密（剩余价值的生产）。前者说明“社会主义代替资本主义是历史必然”，后者说明“资本主义为什么必然被代替、靠什么力量被代替”。两块基石一立，社会主义就不再是道德愿望，而成了科学结论。

\subhdr{2.6\quad 鲜明特征：五个词各管一句话}

\begin{itemize}[leftmargin=1.9em]
\item \dw{科学性}：以客观事实为依据，以事物规律为对象，以实践为检验标准。
\item \dw{革命性}：不满足于解释世界，目的在于改造世界；批判旧世界、建设新世界。
\item \dw{实践性}：在实践中产生、在实践中检验、随实践发展；这是马克思主义区别于其他理论的一个显著特征。
\item \dw{人民性}：第一次站在人民立场上，指明人民群众是历史的创造者，以实现人的自由全面发展为最高价值。
\item \dw{发展性}：与时俱进，随时代、实践和科学的发展而不断丰富。
\end{itemize}

\subhdr{2.7\quad 小例：四句话，判断归属}

下面四句话是这一节最常见的四种错法，逐句给出判断依据：

\begin{enumerate}[label=（\arabic*）,leftmargin=2.4em]
\item “空想社会主义是马克思主义的三个组成部分之一。”——\dw{错}。空想社会主义是\dw{来源}；组成部分是科学社会主义。
\item “德国古典哲学是马克思主义哲学的直接理论来源。”——\dw{对}。黑格尔辩证法与费尔巴哈唯物主义被批判继承，成为马克思主义哲学的直接思想材料。
\item “英国古典政治经济学提供了劳动价值论，是马克思主义政治经济学的来源。”——\dw{对}。斯密、李嘉图的劳动价值论被改造为剩余价值理论的基础。
\item “科学社会主义是马克思主义的理论来源。”——\dw{错}。科学社会主义是\dw{组成部分}（而且是核心与结论），不是来源；它的来源是英法空想社会主义。
\end{enumerate}

\subhdr{2.8\quad 本节错项清单}

\begin{center}\small
\begin{tabular}{@{}p{6.0cm}p{9.4cm}@{}}
\truletop
\thead{错误说法} & \thead{为什么错／正确表述}\\
\trulemid
把空想社会主义列为“组成部分” & 它是理论来源；组成部分是科学社会主义\\
\trulemid
把“马克思主义哲学”说成理论来源 & 它是组成部分（理论基础）；来源是德国古典哲学\\
\trulemid
“《共产党宣言》标志着马克思主义最终确立” & 教材表述是“公开问世”；“最终确立”不是这一句的定位词\\
\trulemid
“马克思主义的三个来源是哲学、政治经济学、科学社会主义” & 这是\dw{组成部分}；三个来源是德国古典哲学、英国古典政治经济学、英法空想社会主义\\
\trulebot
\end{tabular}
\end{center}

\hdr{§3}{第一章：三大规律的“定位词”为什么不能互换}

这是你原卷第 4 题失分的地方，也是整套卷子里最容易补回来的一处：三大规律的内容你可能都背过，但教材给每条规律配了一句\dw{定位词}，考试问“揭示”时，考的就是这句定位词。先把定位词钉死，再回头补内容。

\subhdr{3.1\quad 先立框架：唯物辩证法由什么组成}

唯物辩证法是一个体系，不是三条并列的道理。它的结构是：\dw{总特征}（联系的观点、发展的观点）+\dw{三大规律}（对立统一、质量互变、否定之否定）+\dw{五对范畴}（内容与形式、本质与现象、原因与结果、必然与偶然、现实与可能）+\dw{辩证思维方法}（归纳与演绎、分析与综合、抽象与具体、逻辑与历史相统一）。

三大规律回答的是三个不同的问题，而这三个问题恰好对应教材的三句定位词：

\begin{center}\small
\begin{tabular}{@{}p{3.0cm}p{3.4cm}p{3.2cm}p{4.6cm}@{}}
\truletop
\thead{规律} & \thead{回答的问题} & \thead{教材定位词} & \thead{一句话理解}\\
\trulemid
对立统一规律 & 事物\textbf{为什么}发展 & \dw{源泉和动力} & 内部矛盾推动事物自己运动\\
\trulemid
质量互变规律 & 事物\textbf{怎样}发展 & \dw{形式和状态} & 量变积累到关节点引起质变\\
\trulemid
否定之否定规律 & 事物\textbf{向何处}发展 & \dw{方向和道路} & 螺旋式上升、波浪式前进\\
\trulebot
\end{tabular}
\end{center}

这三句是教材第一章第三节的固定表述。考场上的动作很简单：看到“揭示了事物发展的……”，先在“动力／状态／方向”三项里定位，再看选项把哪一项安到了哪条规律上。

\subhdr{3.2\quad 对立统一规律：为什么它是辩证法的实质和核心}

\textbf{内容。}矛盾是事物内部诸要素之间以及事物之间的对立和统一。同一性指矛盾双方相互依存、相互贯通，在一定条件下相互转化；斗争性指相互排斥、相互分离。同一性是相对的、有条件的，斗争性是绝对的、无条件的，二者结合构成事物的矛盾运动。矛盾推动事物发展，内因是根据，外因是条件，外因通过内因起作用。

\textbf{定位词。}教材说它\dw{揭示了事物发展的源泉和动力}。这句话的含义是：发展不是外力推动的，而是事物内部矛盾双方又统一又斗争的结果。

\textbf{为什么是实质和核心。}三条理由要能说出来：第一，它揭示了事物普遍联系的根本内容和变化发展的内在动力；第二，它是贯穿唯物辩证法其他规律和范畴的中心线索（质量互变、否定之否定都要用矛盾来说明）；第三，矛盾分析方法是认识事物的根本方法。正因为如此，\dw{矛盾的观点是唯物辩证法的根本观点}——这句判断属于第一章，不属于第二章。

\textbf{方法论。}承认矛盾、分析矛盾、解决矛盾；坚持两点论与重点论的统一；具体问题具体分析。

\subhdr{3.3\quad 质量互变规律：度、关节点与“形式和状态”}

\textbf{三个概念。}\dw{质}是一事物区别于他事物的内在规定性；\dw{量}是事物存在和发展的规模、程度、速度等可以用数量表示的规定性；\dw{度}是保持事物质的稳定性的数量界限，度的两端叫关节点（临界点）。

\textbf{两种变化。}量变是数量的增减和场所的变更，是事物在度的范围内的延续和渐进；质变是事物根本性质的变化，是由一种质态向另一种质态的飞跃。

\textbf{辩证关系。}量变是质变的必要准备，质变是量变的必然结果；量变和质变相互渗透——总的量变过程中有部分质变，质变过程中有量的扩张。

\textbf{定位词。}它\dw{揭示了事物发展的形式和状态}：发展采取“渐进—飞跃”的形式，表现为“量变—质变”的状态。

\textbf{方法论。}注重量的积累（不积跬步无以至千里）；抓住时机、促成质变；坚持适度原则（过犹不及）。

\subhdr{3.4\quad 否定之否定规律：扬弃、周期与“方向和道路”}

\textbf{辩证否定观。}否定是事物的自我否定（内部矛盾运动的结果），不是外力强加的；否定是发展的环节（促成新事物产生），也是联系的环节（新旧事物相联系）；否定的实质是\dw{扬弃}——既克服又保留。形而上学否定观则把否定看成全盘抛弃或全盘肯定，或者看成外力作用的结果。

\textbf{否定之否定。}事物经过两次否定、三个阶段（肯定—否定—否定之否定），形成一个周期。周期的终点仿佛回到了起点，但是在更高阶段上的回复，因此发展是\dw{螺旋式上升、波浪式前进}，而不是简单的循环。

\textbf{定位词。}它\dw{揭示了事物发展的方向和道路}：方向是前进的、上升的，道路是迂回的、曲折的。教材用“前途是光明的，道路是曲折的”来概括。

\textbf{方法论。}坚持辩证否定观，反对肯定一切或否定一切；对前途充满信心，又准备走曲折的路；正确对待新生事物。

\subhdr{3.5\quad 三者对照表：考场上的“定位词速查”}

\begin{center}\small
\begin{tabular}{@{}p{2.3cm}p{3.1cm}p{2.3cm}p{3.4cm}p{4.4cm}@{}}
\truletop
\thead{规律} & \thead{回答的问题} & \thead{定位词} & \thead{关键词} & \thead{典型表述}\\
\trulemid
对立统一 & 为什么发展 & 源泉和动力 & 矛盾、同一性与斗争性、内外因、两点论与重点论 & 社会基本矛盾是社会发展的根本动力\\
\trulemid
质量互变 & 怎样发展 & 形式和状态 & 质、量、度、关节点、量变与质变 & 不积跬步无以至千里；过犹不及\\
\trulemid
否定之否定 & 向何处发展 & 方向和道路 & 扬弃、两次否定三阶段、螺旋式上升 & 前途是光明的，道路是曲折的\\
\trulebot
\end{tabular}
\end{center}

\subhdr{3.6\quad 三个小例：让每条规律“只对一件事负责”}

\textbf{例 1（否定之否定）：麦粒—麦株—麦粒。}第一粒麦粒长成麦株（否定），麦株结出新麦粒（否定之否定）。新麦粒不是回到原点：数量多了几十倍，品种还可能被改良。所以这是\dw{在更高阶段上的回复}，说明的是发展的方向和道路；如果有人把它讲成“周而复始的循环”，就把“方向和道路”降格成了循环论，正好丢掉这条规律的定位词。

\textbf{例 2（质量互变）：常压下把水从 20 ℃ 加热到 100 ℃。}0—100 ℃ 之间水保持液态，这是它的“度”；温度是量变，水还是水；到 100 ℃（关节点）继续吸热，水沸腾汽化，发生质变。这个例子只说明“形式和状态”：量变在度的范围内进行，突破关节点才有质变。它说明不了“方向和道路”，也说明不了“源泉和动力”。

\textbf{例 3（对立统一）：商品的使用价值与价值。}任何人不能同时占有二者（对立），二者又统一在同一个商品里（统一）；正是这个矛盾推动商品交换、推动商品经济发展。这里出现的是“矛盾推动发展”，所以对应\dw{源泉和动力}。

\subhdr{3.7\quad 你的原错项与三个自问}

原卷第 4 题你选了“否定之否定规律揭示了事物发展的源泉和动力”。这不是内容不会，而是定位词放错了位置。以后遇到这类题，按顺序自问三句：

\begin{enumerate}[label=\arabic*.,leftmargin=2.0em]
\item 题干问的是“为什么发展”“怎样发展”还是“向何处发展”？
\item 这三个问题分别对应源泉动力、形式状态、方向道路，选项把它们对应对了吗？
\item 选项里有没有把两条规律的内容混在一句话里（比如“否定之否定揭示了源泉和动力”）？
\end{enumerate}

\subhdr{3.8\quad 矛盾学说的补全：内容要能支撑“根本观点”}

定位词之外，矛盾学说本身还有几个固定考点，它们同时是“矛盾的观点是根本观点”的依据：

\begin{itemize}[leftmargin=1.9em]
\item \dw{普遍性与特殊性}：矛盾无处不在、无时不有（普遍性）；不同事物的矛盾各有其特点（特殊性）。共性与个性、绝对与相对的关系是\dw{矛盾问题的精髓}。
\item \dw{主要矛盾与次要矛盾}：在矛盾体系中处于支配地位、对事物发展起决定作用的矛盾是主要矛盾。
\item \dw{矛盾的主要方面与次要方面}：在矛盾统一体中处于支配地位、起主导作用的一方是矛盾的主要方面；事物的性质主要由主要矛盾的主要方面决定。
\item \dw{两点论与重点论的统一}：既看到矛盾的两个方面，又抓住主要矛盾和矛盾的主要方面。
\end{itemize}

\subhdr{3.9\quad 联系、发展与五对范畴（简表）}

联系的特点是客观性、普遍性、多样性、条件性；发展的实质是\dw{新事物的产生和旧事物的灭亡}。判断新旧事物不能看时间先后、形式新奇或力量强弱，而要看它是否符合事物发展的必然趋势、是否具有强大生命力和远大前途。

\begin{center}\small
\begin{tabular}{@{}p{2.4cm}p{6.6cm}p{3.4cm}@{}}
\truletop
\thead{范畴} & \thead{含义要点} & \thead{方法论要点}\\
\trulemid
内容与形式 & 内容是事物的一切要素的总和，形式是内容诸要素的结构和表现方式 & 内容决定形式，形式反作用于内容\\
\trulemid
本质与现象 & 本质是事物的根本性质，现象是本质的外在表现（真象与假象） & 透过现象看本质\\
\trulemid
原因与结果 & 引起与被引起的关系；因果联系具有客观性、复杂性 & 分析因果，预见后果\\
\trulemid
必然与偶然 & 必然性由内部矛盾决定，偶然性由外部条件造成 & 把握必然，重视偶然\\
\trulemid
现实与可能 & 现实是已经实现的可能，可能是尚未实现的现实（区分现实可能与抽象可能） & 立足现实，争取好的可能\\
\trulebot
\end{tabular}
\end{center}

\warnbox{\textbf{本节排雷}\quad 见到“否定之否定揭示源泉和动力”“质量互变揭示方向和道路”“对立统一揭示形式和状态”，直接排除。见到“矛盾的观点是辩证唯物主义的首要的基本观点”，也要停一下：矛盾观点属于\dw{辩证法}，认识论的首要基本观点是\dw{实践}（§4.1 详解）。}

% ======================= PART3 =======================

\hdr{§4}{第二章：实践观点与“能动反映论”}

你在这一章丢了两处：第 6 题把“矛盾的观点”当成了认识论的首要基本观点，多选第 3 题误选了“对事物本来面目做直观反映”。这两处其实是同一件事的两面——\dw{没有把“认识论”这一层与“辩证法”那一层分开，也没有把马克思主义的反映论与旧唯物主义的反映论分开}。

\subhdr{4.1\quad 两个“基本观点”各归其层}

教材里有两个很容易串门的“根本／首要”，它们不在同一章，也不回答同一个问题：

\begin{center}\small
\begin{tabular}{@{}p{5.4cm}p{2.6cm}p{4.0cm}p{3.4cm}@{}}
\truletop
\thead{命题} & \thead{所属层次} & \thead{教材位置} & \thead{为什么}\\
\trulemid
矛盾的观点是唯物辩证法的\textbf{根本观点} & 辩证法 & 第一章第三节 & 对立统一规律是辩证法的实质和核心\\
\trulemid
实践的观点是辩证唯物主义认识论\textbf{首要的和基本的观点} & 认识论 & 第二章第一节 & 实践是认识的基础：来源、动力、目的、检验标准\\
\trulebot
\end{tabular}
\end{center}

第 6 题的题干问的是“认识论”的首要基本观点，正确项自然是\dw{实践}。你选“矛盾的观点”并非全无道理——它确实是辩证法里最根本的观点——但层级不对：**在辩证法里最根本，不等于在认识论里首要**。判断这类题的顺序是：先看题干把问题放在哪一章，再看那一章自己的“第一观点”是谁。

\subhdr{4.2\quad 实践：本质、特征、结构与形式}

\textbf{本质。}实践是人类\dw{能动地改造世界}的社会性的客观物质活动。三个要点缺一不可：它是“能动”的（有目的、有意识），是“改造世界”的（不是单纯的认识或静观），是“社会的、客观物质的”活动（受客观条件制约，具有社会历史性）。

\textbf{三个基本特征。}\dw{直接现实性}——实践能把理论、目的变成现实的存在，这是它区别于认识的根本特点；\dw{自觉能动性}——实践是有目的、有意识的活动；\dw{社会历史性}——实践受社会历史条件制约，并随历史发展而变化。

\textbf{基本结构。}实践主体（具有实践能力的人）、实践客体（进入实践范围的对象）、实践中介（工具、手段以及运用它们的方式方法）。

\textbf{基本形式。}\dw{物质生产实践}（最基本的形式，决定其他一切实践活动）、\dw{社会政治实践}、\dw{科学文化实践}。随着信息技术发展，还出现了虚拟实践这一新形式（教材作为新变化提及，属于拓展，不是基本形式之一）。

\textbf{实践对认识的决定作用（四条）。}实践是认识的\dw{来源}（认识在实践中产生）；实践是认识发展的\dw{动力}（提出新问题、提供新手段）；实践是认识的\dw{目的}（认识为实践服务）；实践是\dw{检验认识真理性的唯一标准}。四条要能一口气说出来——多选题常在这里拆条。

\subhdr{4.3\quad 三条认识路线：先分清“从物到感觉”还是“从感觉到物”}

\begin{center}\small
\begin{tabular}{@{}p{3.4cm}p{4.6cm}p{7.4cm}@{}}
\truletop
\thead{路线} & \thead{核心主张} & \thead{代表与问题}\\
\trulemid
唯物主义反映论 & 从物到感觉和思想：认识是主体对客体的反映 & 坚持反映论，但旧唯物主义把反映理解为消极被动的直观\\
\trulemid
唯心主义先验论 & 从思想和感觉到物：认识先于经验、先于实践 & “生而知之”“天赋观念”“理念回忆”；把认识看成主观自生的东西\\
\trulemid
不可知论 & 世界不可认识，或不能彻底认识 & 休谟、康德；马克思主义用\dw{实践}驳倒不可知论——实践能够把“自在之物”变成“为我之物”\\
\trulebot
\end{tabular}
\end{center}

\subhdr{4.4\quad 能动反映论与直观反映论：这次最容易错的一处}

两条路线都叫“反映论”，差别在\dw{怎么反映}：

\begin{center}\small
\begin{tabular}{@{}p{3.4cm}p{6.2cm}p{6.0cm}@{}}
\truletop
\thead{比较项} & \thead{旧唯物主义（形而上学唯物主义）} & \thead{辩证唯物主义}\\
\trulemid
反映方式 & 消极、被动、直观的反映，像照镜子 & 在实践基础上的\dw{能动}反映\\
\trulemid
与实践的关系 & 脱离实践，把认识看成静观 & 实践是认识的基础，认识随实践发展\\
\trulemid
与辩证法的关系 & 脱离辩证法，不懂得认识的矛盾运动 & 把辩证法应用于认识过程（两次飞跃、反复无限）\\
\trulemid
主体性 & 看不到主体的选择、建构与创造 & 承认认识具有\dw{摹写性}与\dw{创造性}：既反映对象，又能加工、重构、超前反映\\
\trulebot
\end{tabular}
\end{center}

\quoteorig{“最蹩脚的建筑师从一开始就比最灵巧的蜜蜂高明的地方，是他在用蜂蜡建筑蜂房以前，已经在自己的头脑中把它建成了。”\hfill ——马克思《资本论》第一卷}

这句话说的就是“能动”：人的反映不是被动接收，而是在观念中先有目的、有设计，再通过实践把它实现出来。反过来，“直观反映”“照镜子式的反映”“消极被动的反映”这些说法，在教材里是\dw{旧唯物主义}的标签；把它们安到马克思主义认识论头上，就是多选第 3 题 D 项的错法。

\subhdr{4.5\quad 小例：设计路径类材料该怎么读}

材料大意：设计者不预先画出全部道路，而是让人先自由行走，再按人们踩出的痕迹铺设路径。

\textbf{直观反映的读法。}把认识看成对现成对象的被动复写：对象“本来就在那里”，设计者只需“看”和“照抄”。这种读法抹掉了材料里真正发生的事——设计者先有目的（路径要服务于人的行走），再通过实践获得材料（观察足迹），然后比较、概括、形成方案，最后铺路并继续调整。

\textbf{能动反映的读法。}认识在实践基础上展开，走的是“实践—认识—再实践”的路：目的是实践提出的，材料是实践提供的，方案要回到实践中检验。其中既有对对象的摹写（足迹的分布是客观的），也有主体的创造（把足迹概括为路径方案，这是观念中的重构）。这正是教材所说的“主体在实践基础上对客体的能动反映”。

\textbf{顺带说明。}在这类材料里，“事物的本来面目”本身就是人的实践产物（路径是人走出来的），所以 D 项不仅贴错了标签，前提也站不住。多选题里遇到“直观反映”“照镜子”“消极被动”这类词，先排除，再回看其余选项。

\subhdr{4.6\quad 认识的辩证运动：两次飞跃与总规律}

\textbf{感性认识与理性认识。}感性认识包括感觉、知觉、表象，是对事物现象、外部联系的反映；理性认识包括概念、判断、推理，是对事物本质、内部联系的反映。二者是同一认识过程的两个阶段，理性认识依赖于感性认识，感性认识有待于发展为理性认识。

\textbf{第一次飞跃：感性认识 $\rightarrow$ 理性认识。}条件是：感性材料十分丰富且合乎实际；运用科学的思维方法（归纳与演绎、分析与综合、抽象与具体）进行加工。

\textbf{第二次飞跃：理性认识 $\rightarrow$ 实践。}这次飞跃\dw{更重要}：认识世界的目的是改造世界；而且理性认识只有回到实践中才能得到检验、补充和发展。

\textbf{反复性与无限性。}由于客观事物是复杂的、实践是发展的、主体的认识能力受条件限制，一个正确认识往往要经过实践、认识、再实践、再认识的多次反复。认识运动的总规律是：从实践到认识、从认识到实践，循环往复以至无穷，而每一循环都进到高一级的程度。

\subhdr{4.7\quad 真理与价值：本章的第二个考点群}

\textbf{真理的客观性。}真理是人们对客观事物及其规律的正确反映。真理的内容是客观的，检验真理的标准（实践）也是客观的。\dw{真理的绝对性}指任何真理都包含不依赖于人的客观内容，人类能够正确认识无限发展的世界；\dw{真理的相对性}指任何真理在广度上有待扩展、在深度上有待深化。

\textbf{真理与谬误。}相比较而存在、相斗争而发展；在一定条件下相互转化（超出适用范围，真理会变成谬误）。

\textbf{实践是检验真理的唯一标准。}这由真理的本性（主观认识与客观对象相符合）和实践的特点（直接现实性）决定。实践标准是\dw{确定的}（最终能检验）又是\dw{不确定的}（一定历史阶段的实践有局限，检验是一个过程）。

\textbf{价值。}哲学上的价值指客体对主体需要的意义。价值具有\dw{主体性}（因人而异）、\dw{客观性}（以客体的客观属性为基础）、\dw{多维性}、\dw{社会历史性}。价值评价以主客体的价值关系为对象，其正确与否依赖于对客体和主体的正确认识，最终也要由实践来检验。

\textbf{真理与价值的统一。}真理是价值的前提（不符合规律的东西不可能真正满足人的需要），价值是真理的目的（认识真理为了改造世界、满足需要）；二者统一于实践。

\subhdr{4.8\quad 本节排雷词表}

\begin{center}\small
\begin{tabular}{@{}p{4.6cm}p{10.6cm}@{}}
\truletop
\thead{见到就排除的说法} & \thead{正确表述或理由}\\
\trulemid
“认识是对事物本来面目的直观反映” & 旧唯物主义的直观反映论；马克思主义是能动的反映论\\
\trulemid
“认识是主体主观自生的” & 唯心主义先验论；认识来源于实践\\
\trulemid
“感性认识比理性认识更可靠” & 二者各有局限；理性认识依赖感性认识，感性认识有待上升为理性认识\\
\trulemid
“有用就是真理” & 实用主义真理观；真理的检验标准是实践，不是效用\\
\trulemid
“真理是多元的” & 真理具有客观性，对同一对象的真理性认识只有一个\\
\trulemid
“认识可以一次完成” & 认识具有反复性和无限性\\
\trulebot
\end{tabular}
\end{center}

% ======================= PART4 =======================

\hdr{§5}{第三章：唯物史观的骨架，与“革命的根本问题”}

这一章你做得最好——社会存在与社会意识、生产方式、经济基础与上层建筑、资本主义再生产实质都一次做对。所以本节不再从头讲，而是把骨架压缩成可复述的条目，重点补上你错的那一处：\dw{无产阶级革命的根本问题为什么是国家政权}。

\subhdr{5.1\quad 社会存在与社会意识：三要素、两句话}

\textbf{社会存在}是社会生活的物质方面，包括三个要素：\dw{自然地理环境}、\dw{人口因素}、\dw{物质生产方式}。前两者是社会发展的必要条件，能加速或延缓社会发展，但不能决定社会性质；\dw{物质生产方式}才是社会历史发展的决定力量——它决定社会的性质、结构和面貌，决定社会形态的更替。

\textbf{社会意识}是社会生活的精神方面，是社会存在的反映。按水平分，有社会心理（自发、感性）与社会意识形式（自觉、系统）；按是否反映经济基础分，有意识形态（政治法律思想、道德、艺术、宗教、哲学等）与非意识形态（自然科学、语言学、逻辑学等）。

\textbf{两句话。}第一，\dw{社会存在决定社会意识}：社会意识的内容来源于社会存在，随社会存在的变化而变化。第二，\dw{社会意识具有相对独立性}：它与社会存在的发展不完全同步（可能超前或落后）、不平衡（经济落后国家的某些精神领域可能领先）、有历史继承性、各形式之间相互影响，并且\dw{对社会存在具有能动的反作用}——先进的社会意识促进社会发展，落后的社会意识阻碍社会发展。

\textbf{与唯心史观的分界。}唯心史观有两个主要缺陷：只看人们的思想动机，不追究动机背后的物质根源；否认人民群众的历史作用。唯物史观的回答是：社会存在决定社会意识，人民群众是历史的创造者。

\subhdr{5.2\quad 两大规律：生产力与生产关系、经济基础与上层建筑}

\begin{center}\small
\begin{tabular}{@{}p{3.0cm}p{6.2cm}p{6.4cm}@{}}
\truletop
\thead{规律} & \thead{内容要点} & \thead{结论与意义}\\
\trulemid
生产力与生产关系 & 生产力三要素：劳动者（最活跃）、劳动资料（生产工具是重要标志）、劳动对象；科学技术是第一生产力。生产关系三方面：生产资料所有制（最基本、决定性的）、生产中人与人的关系、产品分配关系 & \dw{生产关系一定要适合生产力状况}：生产力决定生产关系，生产关系反作用于生产力\\
\trulemid
经济基础与上层建筑 & 经济基础是占统治地位的生产关系的总和；上层建筑包括意识形态与政治法律制度及设施，\dw{国家政权是核心} & \dw{上层建筑一定要适合经济基础状况}：经济基础决定上层建筑，上层建筑反作用于经济基础\\
\trulebot
\end{tabular}
\end{center}

\textbf{国家。}国家是阶级矛盾不可调和的产物，是经济上占统治地位的阶级进行阶级统治的工具。\dw{国体}指国家的阶级性质（哪个阶级掌握政权），\dw{政体}指政权的组织形式。这个“工具”的定性，是后面理解“革命的根本问题是国家政权问题”的前提。

\subhdr{5.3\quad 社会形态更替：三条统一}

社会形态是经济基础与上层建筑的统一体。社会形态更替有三个“统一”：\dw{统一性与多样性}的统一（五种社会形态依次更替是总趋势，但不同国家可以跨越或缩短某些阶段）；\dw{必然性与选择性}的统一（更替有客观规律，但主体的历史选择可以加速或延缓）；\dw{前进性与曲折性}的统一（总方向是前进的，具体道路可能反复）。

\subhdr{5.4\quad 动力系统与“革命的根本问题”}

社会发展的动力是一个系统，按层次记：

\begin{itemize}[leftmargin=1.9em]
\item \dw{社会基本矛盾}（生产力与生产关系、经济基础与上层建筑的矛盾）是社会发展的\dw{根本动力}。
\item \dw{阶级斗争}是阶级社会发展的\dw{直接动力}；阶级的产生、阶级斗争的根源都在于经济关系。
\item \dw{社会革命}是社会形态的更替，是革命阶级推翻反动阶级的统治、用新的社会制度代替旧的社会制度，是“历史的火车头”。
\item \dw{改革}是同一社会形态内部的量变和部分质变；社会主义改革是社会主义制度的自我完善和发展。
\item \dw{科学技术}是推动社会发展的强大杠杆，但它是双刃剑，其社会后果取决于社会制度与使用方式。
\end{itemize}

\textbf{革命的“根本问题”。}教材在讲社会革命时给出这一判断：\dw{革命的根本问题是国家政权问题}。列宁在《论两个政权》中也说“一切革命的根本问题是国家政权问题”。道理不难推：革命要改变旧的生产关系，而旧的生产关系正是靠旧的国家政权来维护的；不夺取政权，就不能改变生产关系，也就谈不上社会形态的更替。

\textbf{为什么不是“统一战线中的领导权问题”。}统一战线的领导权，是无产阶级在革命同盟军中的关键问题（中国革命的经验把它概括为“三大法宝”之一相关联的问题），它重要，但它是\dw{策略与领导权}层面的问题；题干问的是“一切革命的根本问题”，答案只能落在\dw{政权}上——因为没有政权，其他一切都无从实现。你原来选 A，是把“重要问题”误当成了“根本问题”。

\textbf{一章归属的说明。}教材在第三章第二节讲社会革命时给出“革命的根本问题”，到第六章讲社会主义代替资本主义时，结论以“无产阶级革命的根本问题是国家政权问题”的形式再次出现。所以无论你的课件把它放在第三章还是第六章，结论相同；考场上按“国家政权”作答。

\subhdr{5.5\quad 群众史观：三条创造、两条路线}

\textbf{人民群众是历史的创造者}，体现在三个方面：是社会物质财富的创造者（生产力主体），是社会精神财富的创造者（实践提供源泉），是社会变革的决定力量。

\textbf{群众观点与群众路线。}群众观点：相信人民群众自己解放自己，全心全意为人民服务，一切向人民群众负责，虚心向人民群众学习。群众路线：一切为了群众，一切依靠群众，从群众中来，到群众中去。

\textbf{评价历史人物。}坚持\dw{历史分析法}（放在当时的历史条件下）与\dw{阶级分析法}（分析其阶级属性），反对英雄史观，也反对否定个人的历史作用。

\subhdr{5.6\quad 本节排雷}

\begin{center}\small
\begin{tabular}{@{}p{5.2cm}p{10.0cm}@{}}
\truletop
\thead{错误说法} & \thead{正确表述}\\
\trulemid
“地理环境决定社会性质” & 地理环境是必要条件；决定力量是物质生产方式\\
\trulemid
“社会意识总是落后于社会存在” & 相对独立性包括不完全同步：可能超前，也可能落后\\
\trulemid
“上层建筑决定经济基础” & 经济基础决定上层建筑，上层建筑反作用于经济基础\\
\trulemid
“革命的根本问题是统一战线领导权” & 革命的根本问题是国家政权问题\\
\trulemid
“改革也是革命，因为要改变根本制度” & 改革是同一社会形态内部的量变和部分质变，社会主义改革是制度的自我完善\\
\trulebot
\end{tabular}
\end{center}

% ======================= PART5 =======================

\hdr{§6}{第四章：资本主义的本质——从劳动力商品到剩余价值}

这一章你丢了两处：\dw{劳动力商品价值的三要素}只记住第一条（多选第 4 题），以及\dw{货币转化为资本的前提}最初留空（第 10 题）。前者是条目残缺，后者是因果链断裂。本章按“商品—货币—资本—剩余价值—积累—循环周转—再生产—分配”的顺序走一遍，每个环节都给公式和最简数字。

\subhdr{6.1\quad 商品二因素与劳动二重性（政治经济学的枢纽）}

\begin{center}\small
\begin{tabular}{@{}p{2.6cm}p{4.4cm}p{4.4cm}p{4.0cm}@{}}
\truletop
\thead{项目} & \thead{使用价值} & \thead{价值} & \thead{关系}\\
\trulemid
属性 & 自然属性：能满足人们某种需要 & 社会属性：凝结在商品中的无差别人类劳动 & 二者对立统一于商品之中\\
\trulemid
由什么创造 & 具体劳动（有用劳动） & 抽象劳动（无差别的人类劳动） & 劳动二重性决定商品二因素\\
\trulemid
能否同时占有 & 卖出后生产者失去使用价值 & 卖出后生产者实现价值 & 任何人不能同时占有二者\\
\trulebot
\end{tabular}
\end{center}

\textbf{劳动二重性。}具体劳动创造使用价值，并在生产中把生产资料的旧价值转移到新产品中；抽象劳动形成价值。教材把劳动二重性学说称为\dw{理解马克思主义政治经济学的枢纽}——因为只有它才能说明价值由什么创造、剩余价值从哪里来。

\subhdr{6.2\quad 价值量、货币与价值规律}

\textbf{价值量。}商品的价值量由生产商品的\dw{社会必要劳动时间}决定：在现有的社会正常的生产条件下，在社会平均的劳动熟练程度和劳动强度下制造某种使用价值所需要的劳动时间。劳动生产率与单位商品的价值量成反比：社会劳动生产率提高，单位商品价值量下降。

\textbf{货币。}货币是价值形式长期发展的结果，是固定地充当一般等价物的商品。它有价值尺度、流通手段、贮藏手段、支付手段、世界货币五种职能，其中价值尺度和流通手段是基本职能。

\textbf{价值规律。}内容是：商品的价值量由社会必要劳动时间决定，商品交换以价值量为基础实行等价交换。表现形式是价格围绕价值上下波动（这不否定价值规律，因为波动的中心仍是价值）。作用是：自发调节生产资料和劳动力在社会各生产部门之间的分配；刺激商品生产者改进技术、提高劳动生产率；调节商品生产者的收入分配，导致优胜劣汰。

\textbf{私有制商品经济的基本矛盾。}\dw{私人劳动与社会劳动的矛盾}：生产者的劳动既是私人劳动（私人支配），又必须是社会劳动（产品要卖出去）。这个矛盾决定商品生产者的命运，是私有制商品经济一切矛盾的根源。

\subhdr{6.3\quad 货币转化为资本：一步不能少的因果链}

\textbf{资本总公式。}G—W—G$'$，其中 G$'$ = G + $\Delta$G。货币转化为资本，关键在于 $\Delta$G（剩余价值）从哪里来。

\textbf{矛盾。}按价值规律，等价交换不产生价值，所以剩余价值不能在流通中产生；但如果不进入流通、不买不卖，价值也不会增殖，所以剩余价值又不能离开流通。

\textbf{解决。}答案在一种特殊商品上——\dw{劳动力}。劳动力商品的使用价值有一个别的商品没有的特点：它在使用中\dw{能创造出比自身价值更大的价值}。所以货币所有者买到劳动力之后，使用它就能获得剩余价值，货币也就转化为资本。

这条链要能一口气说出来：\dw{劳动力成为商品 $\rightarrow$ 货币转化为资本 $\rightarrow$ 剩余价值生产 $\rightarrow$ 资本主义生产方式}。第 10 题问的“前提”，正是链条的第一环。

\textbf{劳动力成为商品的两个基本条件。}第一，劳动者有人身自由，能够把自己的劳动力当作自己的商品来支配；第二，劳动者丧失一切生产资料和生活资料，“自由得一无所有”，除了劳动力之外没有别的商品可以出卖。两个条件缺一不可：只有自由而没有“一无所有”，劳动者可以自营，不必出卖劳动力；只有一无所有而没有人身自由，就是奴隶，劳动力不属于自己。

\subhdr{6.4\quad 劳动力商品的价值：三要素必须完整}

\textbf{决定原则。}劳动力商品的价值，由生产、发展、维持和延续劳动力所必需的生活必需品的价值决定。它包括三个部分：

\begin{enumerate}[label=（\arabic*）,leftmargin=2.4em]
\item \dw{维持劳动者本人生存所必需的生活资料的价值}——人要活着才能劳动，这是最低的一条。
\item \dw{维持劳动者家属（家庭）生存所必需的生活资料的价值}——劳动力会衰老、死亡，要有下一代接续，劳动力再生产是代际的再生产。
\item \dw{劳动者接受教育和训练所支出的费用}——现代生产需要技能和知识，而这些不是天生的，要靠教育和训练获得。
\end{enumerate}

此外，劳动力价值还包含\dw{历史和道德的因素}：不同国家、不同历史时期，生活必需品的内容和水平不同（这就是为什么各国劳动力价值不能直接按同一标准比较）。

\textbf{记忆锚：活着—繁衍—会干。}三项分别对应“本人活着”“下一代接续”“有技能可干”。多选第 4 题你只选了第一项，漏掉了“繁衍”和“会干”，本质上是把“劳动力价值”读成了“本人的生活费”。

\textbf{为什么教育训练费必须算进去。}因为劳动力是存在于活的人体中的能力，技能水平决定这种能力的质量；培养技能的费用是劳动力再生产所必需的支出，构成劳动力价值的一部分。它不是资本家的“额外恩赐”——资本家支付这部分价值，恰恰是因为他需要更高质量的劳动力。

\textbf{小例（数字从简）。}某地一名工人的劳动力价值按月计：本人生活资料 3000 元，家属生活资料 1200 元，教育培训费 300 元，合计 4500 元。如果只算第一项（3000 元），就会低估劳动力价值，进而在计算剩余价值率时高估剥削程度。三项一个都不能少，不只是为了背条目，也是为了让后面的计算站得住。

\textbf{把三要素接到剩余价值率上。}设日劳动力价值 100 元，必要劳动 4 小时；工作日 8 小时，创造价值 200 元。这时 $m=100$ 元，$v=100$ 元，

\[
m'=\frac{m}{v}=\frac{100}{100}=100\%.
\]

若教育培训费上升使劳动力价值提高到 120 元，在工作日长度不变、社会劳动生产率不变的条件下，必要劳动延长到 4.8 小时，剩余劳动缩短为 3.2 小时，于是 $m'=3.2/4.8\approx 66.7\%$。这说明：\dw{工作日长度不变时，劳动力价值与剩余价值此消彼长}——这正是“相对剩余价值”生产要拼命提高劳动生产率、压低劳动力价值的原因。

\subhdr{6.5\quad 剩余价值的生产：两种基本方法}

\textbf{劳动过程与价值增殖过程。}劳动过程的三要素是劳动者、劳动对象、劳动资料。价值增殖过程是“超过一定点而延长了的价值形成过程”：这个“一定点”，就是补偿劳动力价值所需的必要劳动时间。

\textbf{不变资本与可变资本。}用于购买生产资料的资本是\dw{不变资本 c}，它在生产中只转移价值、不改变价值量；用于购买劳动力的资本是\dw{可变资本 v}，它能创造出大于自身价值的价值。区分二者的依据是\dw{资本不同部分在剩余价值生产中的作用}，意义在于揭示剩余价值的唯一源泉是可变资本（即雇佣劳动）。

\textbf{剩余价值率。}剩余价值与可变资本的比率：

\[
m'=\frac{m}{v}=\frac{\text{剩余劳动时间}}{\text{必要劳动时间}}.
\]

\textbf{两种基本方法。}\dw{绝对剩余价值}：必要劳动时间不变，延长工作日（或提高劳动强度）。\dw{相对剩余价值}：工作日长度不变，通过提高全社会劳动生产率缩短必要劳动时间。\dw{超额剩余价值}是相对剩余价值的特殊形式：个别企业率先提高劳动生产率，使商品的个别价值低于社会价值，按社会价值出售而多得的那部分剩余价值。

\textbf{数字例。}工作日 8 小时，必要劳动 4 小时、剩余劳动 4 小时，$m'=4/4=100\%$。

\begin{itemize}[leftmargin=1.9em]
\item 绝对剩余价值：工作日延长到 12 小时，必要劳动仍 4 小时，剩余劳动 8 小时，$m'=8/4=200\%$。
\item 相对剩余价值：工作日仍 8 小时，全社会劳动生产率提高使必要劳动缩短到 2 小时，剩余劳动 6 小时，$m'=6/2=300\%$。
\item 超额剩余价值：某企业生产一件商品的社会价值 10 元、个别价值 8 元，产量 100 件，则超额剩余价值 $=(10-8)\times100=200$ 元。
\end{itemize}

\textbf{生产自动化条件下的剩余价值源泉。}机器不创造价值，价值仍然由活劳动创造。自动化提高了劳动生产率、改变了剩余价值生产的方法（更多依靠相对剩余价值），但没有改变剩余价值的源泉。这个判断常以“机器人取代工人后剩余价值来自哪里”的形式出现，答案始终是：来自雇佣工人的剩余劳动。

\subhdr{6.6\quad 资本积累：有机构成与相对过剩人口}

\textbf{实质。}资本积累是\dw{剩余价值的资本化}——把剩余价值再转化为资本。资本积累的源泉是剩余价值，不是资本家的“节欲”。

\textbf{资本有机构成。}由资本技术构成决定并反映技术构成变化的资本价值构成，记作 $c:v$。例如 $c=8000$、$v=2000$ 时，有机构成为 $4:1$。

\textbf{相对过剩人口。}随着资本积累和有机构成提高，总资本中可变资本的相对比重下降，对劳动力的需求相对减少，形成相对过剩人口（不是绝对过剩，而是相对于资本增殖需要而言的过剩）。它既是资本主义生产方式的必然产物，又是其存在条件（产业后备军）。

\textbf{一般规律与历史趋势。}资本积累的一般规律是：一端是财富的积累，另一端是贫困的积累。其历史趋势是资本主义基本矛盾（生产社会化与生产资料私人占有）日益尖锐，最终指向“两个必然”。

\subhdr{6.7\quad 资本循环与周转}

\textbf{三个阶段、三种职能形式。}

\begin{center}\small
\begin{tabular}{@{}p{2.6cm}p{4.8cm}p{3.6cm}p{4.8cm}@{}}
\truletop
\thead{阶段} & \thead{内容} & \thead{职能形式} & \thead{作用}\\
\trulemid
购买阶段 & 用货币购买生产资料和劳动力（G—W） & 货币资本 & 为剩余价值生产准备条件\\
\trulemid
生产阶段 & 生产资料与劳动力结合，生产出包含剩余价值的商品 & 生产资本 & 生产剩余价值（决定性阶段）\\
\trulemid
售卖阶段 & 卖出商品，收回货币（W$'$—G$'$） & 商品资本 & 实现剩余价值\\
\trulebot
\end{tabular}
\end{center}

\textbf{连续进行的条件。}三种职能形式在\dw{空间上并存}、在\dw{时间上继起}。缺一不可：只并存不继起会停顿，只继起不并存会断流。

\textbf{资本周转。}周转时间 = 生产时间 + 流通时间。生产资本按价值周转方式分为\dw{固定资本}（价值逐次转移，如机器设备）与\dw{流动资本}（价值一次转移，如原料、燃料、工资）。注意：固定资本与流动资本的划分依据是\dw{价值周转方式}，不是物质形态，也不是在剩余价值生产中的作用（后者是 c 与 v 的划分）。

\textbf{周转速度与年剩余价值率。}加快周转可以提高年剩余价值量。设 $n$ 为一年内可变资本的周转次数，

\[
M'=m'\cdot n .
\]

例如 $m'=100\%$、预付可变资本 100 万元、一年周转 6 次，则年剩余价值量 600 万元，年剩余价值率 $M'=600\%$。

\subhdr{6.8\quad 社会再生产：两大部类与三个基本条件}

社会再生产的核心问题是\dw{社会总产品的实现}，即价值补偿（产品卖出去、收回价值）与实物补偿（买到再生产所需的生产资料和消费资料）。

社会总产品按最终用途分为两大部类：\dw{第 I 部类}（生产资料生产）与\dw{第 II 部类}（消费资料生产）。简单再生产的基本实现条件是：

\[
\mathrm{I}(v+m)=\mathrm{II}c ,
\]

即第 I 部类的可变资本与剩余价值之和，必须等于第 II 部类的不变资本。由此可推出另外两个条件：

\[
\mathrm{I}(c+v+m)=\mathrm{I}c+\mathrm{II}c,\qquad
\mathrm{II}(c+v+m)=\mathrm{I}(v+m)+\mathrm{II}(v+m).
\]

第一个条件是最基本的：它说明两大部类之间必须保持比例关系，生产社会化要求按比例分配社会劳动，而私有制商品经济（以及资本主义）无法自觉保持这种比例——这正是资本主义基本矛盾在再生产中的表现。

\subhdr{6.9\quad 工资与剩余价值的分配}

\textbf{工资的本质。}资本主义工资是\dw{劳动力价值或价格的转化形式}，却在现象上表现为“劳动的价格”。这一表现掩盖了必要劳动与剩余劳动的区别：好像工人得到了全部劳动的报酬，剩余价值就“看不见”了。

\textbf{剩余价值怎样被瓜分。}剩余价值先在各部门之间通过竞争平均化，再以各种形式被不同集团占有：

\begin{center}\small
\begin{tabular}{@{}p{3.4cm}p{3.6cm}p{4.4cm}p{4.0cm}@{}}
\truletop
\thead{形式} & \thead{获得者} & \thead{量的规定} & \thead{来源}\\
\trulemid
产业利润 & 产业资本家 & 成本价格 + 平均利润（生产价格） & 产业工人创造的剩余价值\\
\trulemid
商业利润 & 商业资本家 & 参与平均利润率的形成 & 产业工人创造的剩余价值的一部分\\
\trulemid
借贷利息 & 借贷资本家 & 平均利润的一部分 & 同上（职能资本家让渡）\\
\trulemid
银行利润 & 银行资本家 & 存贷利差等 & 同上\\
\trulemid
资本主义地租 & 土地所有者 & 级差地租、绝对地租 & 农业工人创造的剩余价值\\
\trulebot
\end{tabular}
\end{center}

\textbf{两个比率要分清。}剩余价值率 $m'=m/v$ 反映剥削程度；利润率 $p'=m/(c+v)$ 反映资本的增殖程度。因为分母加了不变资本，$p'<m'$。例如 $c=8000$、$v=2000$、$m=2000$，则 $m'=100\%$，而 $p'=2000/10000=20\%$。

\subhdr{6.10\quad 本节排雷}

\begin{center}\small
\begin{tabular}{@{}p{5.6cm}p{9.6cm}@{}}
\truletop
\thead{错误说法} & \thead{正确表述}\\
\trulemid
“资本家全部资本都带来剩余价值” & 只有可变资本带来剩余价值；不变资本只转移价值\\
\trulemid
“剩余价值在流通中产生” & 剩余价值在生产中产生，在流通中实现\\
\trulemid
“工资是劳动的价值或价格” & 工资是劳动力价值或价格的转化形式\\
\trulemid
“劳动力价值只包括本人生活资料” & 还包括家属生活资料与教育培训费用\\
\trulemid
“资本有机构成提高，剩余价值率必然下降” & 有机构成提高可能提高剩余价值率（相对剩余价值），但会减少对劳动力的需求\\
\trulemid
“固定资本与流动资本的划分依据是物质形态” & 划分依据是价值周转方式\\
\trulebot
\end{tabular}
\end{center}

% ======================= PART6 =======================

\hdr{§7}{第五章：垄断资本主义——“五特征”与“性质未变”}

这一章你丢了两处：多选第 5 题只选出“金融寡头统治”，漏掉“垄断组织起决定作用”和“资本输出具有特别重要意义”；第 12 题把国家垄断资本主义说成“改变了经济的资本主义性质”。前者是条目不全，后者是性质判断错位。这一节把五特征按列宁的原话排成一条线，再给“性质未变”三条论证。

\subhdr{7.1\quad 自由竞争怎样走到垄断}

\textbf{垄断的形成。}自由竞争引起生产集中和资本集中；生产集中发展到一定阶段，必然引起垄断——因为少数大企业彼此势均力敌，继续竞争代价太大，于是转向联合或默契，形成垄断组织。

\textbf{垄断与竞争并存。}垄断没有消除竞争，反而使竞争更加激烈、形式更加复杂：垄断组织内部、垄断组织之间、垄断组织与非垄断企业之间都存在竞争。教材的说法是：垄断是从自由竞争中产生的，垄断又加剧了竞争。

\textbf{垄断利润与垄断价格。}垄断组织凭借垄断地位，获得大大超过平均利润的高额利润，即垄断利润。其手段是垄断价格（垄断高价与垄断低价）。\dw{垄断价格并不否定价值规律}：全社会商品的价值总量与价格总额仍然相等，垄断高价与垄断低价相互抵消；改变的只是价值规律作用的表现形式。

\subhdr{7.2\quad 金融资本与金融寡头}

\textbf{金融资本}是工业垄断资本和银行垄断资本融合在一起形成的垄断资本。银行从普通的中介人变成了万能的垄断者。

\textbf{金融寡头}是掌握金融资本、操纵国民经济和国家政权的少数大资本家。它的统治方式：经济上通过\dw{参与制}（控股、层层控制）实现，政治上通过\dw{个人联合}（政商旋转门、资助选举、控制舆论）实现。

\subhdr{7.3\quad 列宁的五个基本经济特征（必须按顺序背）}

列宁在《帝国主义是资本主义的最高阶段》（1916）中概括了垄断资本主义的五个基本经济特征：

\begin{enumerate}[label=（\arabic*）,leftmargin=2.4em]
\item \dw{生产和资本的集中发展到这样高的程度，以致造成了在经济生活中起决定作用的垄断组织}；
\item \dw{银行资本和工业资本已经融合起来，在这个“金融资本”的基础上形成了金融寡头}；
\item \dw{与商品输出不同的资本输出有了特别重要的意义}；
\item \dw{瓜分世界的资本家国际垄断同盟已经形成}；
\item \dw{最大资本主义大国已把世界上的领土瓜分完毕}。
\end{enumerate}

\textbf{记忆锚。}按“内部—主体—对外—同盟—领土”的顺序串：先是\dw{垄断组织}在经济中起决定作用（这是最基础的一条，其余四条都建立在垄断之上）；再是\dw{金融寡头}的统治；然后是对外扩张从商品输出转为\dw{资本输出}；接着是国际垄断\dw{同盟}瓜分世界；最后是列强把世界\dw{领土}瓜分完毕。

\textbf{为什么第一条最基础。}因为垄断组织在经济生活中起决定作用，是垄断资本主义区别于自由竞争资本主义的根本标志；金融资本、资本输出、国际垄断同盟、领土瓜分都是在这个基础上展开的。多选第 5 题你只选了金融寡头，漏掉第一条和第三条，等于把“基础”和“对外扩张的标志”一起丢了。

\textbf{注意“商品输出”与“资本输出”的区别。}自由竞争阶段以商品输出为主；垄断阶段资本输出（对外投资、贷款）具有特别重要的意义。选项里如果写“商品输出具有特别重要的意义”，那是把阶段搞反了。

\subhdr{7.4\quad 国家垄断资本主义：“性质未变”三条论证}

\textbf{含义。}国家垄断资本主义是国家政权和私人垄断资本融合在一起的垄断资本主义。

\textbf{产生原因。}科技进步和生产社会化程度进一步提高，要求国家出面承担私人资本无力或不愿承担的投资；经济危机的频繁爆发要求国家干预；战争和军备竞赛的刺激；国际竞争的加剧。

\textbf{主要形式。}（1）国家所有并直接经营的企业；（2）国家与私人共有、合营的企业；（3）国家通过多种形式参与私人垄断资本的再生产过程（国家采购、补贴、贷款、税收优惠等）；（4）国家通过宏观调节和微观规制干预经济。

\textbf{实质（三条论证，逐条对应“性质未变”）。}

\begin{enumerate}[label=\arabic*.,leftmargin=2.0em]
\item \dw{所有制没变}：国有企业属于资产阶级国家所有，而国家是资产阶级的总代表；国有化不是“公有化”，它没有改变资本雇佣劳动的关系。
\item \dw{目的没变}：国家干预的目标是维护垄断资本的整体利益和资本主义制度的存续，而不是满足人民需要。国家承担风险、私人获取利润，是“总资本家”的行为方式。
\item \dw{结果没变}：剩余价值规律和资本主义基本矛盾仍在起作用；国家干预只能缓和而不能消除危机，反而可能积累新的矛盾（如债务膨胀、资产泡沫）。
\end{enumerate}

因此教材的结论是：国家垄断资本主义\dw{没有改变垄断资本主义的性质}，它只是\dw{资本主义生产关系的局部调整}，从根本上说仍然服从于垄断资产阶级的整体利益。第 12 题选“改变了经济的资本主义性质”，错在把“局部调整”读成了“性质改变”。

\textbf{小例。}2008 年国际金融危机中，西方国家用财政资金救助“大到不能倒”的私人金融机构：出资的是国家（全体纳税人），得救的是私人垄断资本。这个例子正好说明三条论证——所有制没变（钱来自国家但所有权与利润归私人）、目的没变（保的是金融资本体系）、结果没变（危机缓和了，基本矛盾没有解决）。

\subhdr{7.5\quad 经济全球化：四个表现、两重后果}

\textbf{表现（四个）。}\dw{生产全球化}（国际分工与产业链跨国布局）、\dw{贸易全球化}（世界市场与自由贸易体制）、\dw{金融全球化}（资本跨国流动、金融市场联动）、\dw{企业经营全球化}（跨国公司成为主体）。

\textbf{动因。}科技进步和生产力发展提供了物质条件；跨国公司的发展提供了组织载体；各国市场经济体制的普遍确立提供了制度条件；国际经济组织的建立提供了协调机制。

\textbf{二重后果。}积极方面：促进资源在全球范围内优化配置，推动生产力发展，促进科技文化交流。消极方面：发达国家与发展中国家的差距扩大（发展不平衡）；全球性环境问题；金融风险的跨国传染；国际经济规则由发达国家主导，发展中国家话语权不足。

\subhdr{7.6\quad 当代资本主义的新变化：变了什么，没变什么}

\textbf{变了什么（六个方面）。}生产资料所有制（国家资本所有制、法人资本所有制发展）；劳资关系与分配关系（职工参与决策、股票期权、社会福利制度）；社会阶层结构（资本所有权与控制权分离、中间阶层扩大）；经济调节机制（国家干预、计划化、国际协调）；政治制度（公民权利扩大、政党制度与选举制度变化）；意识形态（多元化、淡化意识形态色彩）。

\textbf{为什么变。}科学技术革命和生产力发展的推动；工人阶级争取自身权利的斗争；社会主义制度的影响；主张改良的政党对资本主义制度的改革。

\textbf{没变什么（实质）。}这些变化是\dw{人类社会发展一般规律和资本主义基本矛盾作用的结果}，是\dw{资本主义生产关系的自我调整}；它\dw{没有改变资本主义制度的本质}，\dw{没有克服资本主义基本矛盾}，\dw{没有改变资本主义必然被社会主义取代的历史总趋势}。

\subhdr{7.7\quad 本节排雷}

\begin{center}\small
\begin{tabular}{@{}p{6.0cm}p{9.2cm}@{}}
\truletop
\thead{错误说法} & \thead{正确表述}\\
\trulemid
“国家垄断资本主义改变了经济的资本主义性质” & 没有改变；它是资本主义生产关系的局部调整\\
\trulemid
“垄断消除了竞争” & 垄断与竞争并存，垄断使竞争更激烈、更复杂\\
\trulemid
“垄断价格否定了价值规律” & 没有否定；改变的是价值规律作用的表现形式\\
\trulemid
“国家垄断资本主义已经克服了资本主义基本矛盾” & 只能缓和，不能克服\\
\trulemid
“经济全球化就是贸易全球化” & 还包括生产、金融、企业经营全球化\\
\trulemid
“当代资本主义的新变化说明资本主义性质已变” & 变化是自我调整，本质、基本矛盾、历史趋势均未改变\\
\trulebot
\end{tabular}
\end{center}

% ======================= PART7 =======================

\hdr{§8}{第六、七章：科学社会主义与共产主义}

这一章你丢了两处：第 13 题的“革命根本问题”已在 §5.4 辨析过；第 15 题把“两个绝不会”的出处记成了《共产党宣言》。这一节把社会主义从空想到科学的三段历程压缩成时间线，再把“两个必然”与“两个绝不会”放在一起讲清它们为什么不矛盾。

\subhdr{8.1\quad 空想社会主义的三个阶段}

\begin{center}\small
\begin{tabular}{@{}p{3.4cm}p{3.2cm}p{4.4cm}p{4.2cm}@{}}
\truletop
\thead{阶段} & \thead{时间} & \thead{代表人物与著作} & \thead{特点}\\
\trulemid
早期空想社会主义 & 16—17 世纪 & 莫尔《乌托邦》、康帕内拉《太阳城》 & 用文学游记描绘理想社会，批判私有制\\
\trulemid
空想平均共产主义 & 18 世纪 & 摩莱里、马布利、巴贝夫 & 主张平均主义、禁欲主义，带有小生产色彩\\
\trulemid
批判的空想社会主义 & 19 世纪初 & 圣西门、傅立叶、欧文 & 批判资本主义制度，设想未来社会，但找不到实现的社会力量和道路\\
\trulebot
\end{tabular}
\end{center}

空想社会主义的历史贡献是深刻批判了资本主义，提出了许多合理设想；局限是历史观唯心主义（寄望于天才人物和统治者的善意），因而“空想”。这也正是它只能作为\dw{理论来源}、不能作为\dw{组成部分}的原因（见 §2.4）。

\subhdr{8.2\quad 科学社会主义的创立：两块基石}

唯物史观揭示了社会发展的客观规律，剩余价值学说揭示了资本主义剥削的秘密。两块基石一立，社会主义就从“应该如此”的道德呼吁，变成“必然如此”的科学结论。1848 年《共产党宣言》的发表，标志着科学社会主义的公开问世。

\textbf{社会主义从理论到实践的三次飞跃。}从空想到科学（19 世纪 40 年代）；从理论到现实（1917 年十月革命建立第一个社会主义国家）；从一国到多国（二战后社会主义阵营形成）。中国特色社会主义是科学社会主义在当代中国的实践形态。

\subhdr{8.3\quad 社会主义的一般原则与经济文化落后国家的长期性}

\textbf{一般原则（记核心几条，不背“共几条”）。}以唯物史观为理论基础，坚信资本主义必然灭亡、社会主义必然胜利；坚持无产阶级专政；坚持生产资料公有制和按劳分配；坚持无产阶级政党的领导；以人的自由全面发展为最高价值目标；合乎自然规律地改造和利用自然，实现人与自然和谐共生。

\textbf{经济文化相对落后国家建设社会主义的长期性。}原因有四个方面：生产力基础薄弱，商品经济不发达；面临发达资本主义国家的压力与封锁；没有现成经验可循，只能在实践中探索；旧的思想观念和习惯势力还有影响。所以社会主义在这些国家的发展必然是长期的、艰巨的、曲折的，不能急于求成。

\subhdr{8.4\quad 共产主义社会的基本特征}

教材把共产主义社会的基本特征概括为三条：

\begin{enumerate}[label=（\arabic*）,leftmargin=2.4em]
\item \dw{物质财富极大丰富，消费资料按需分配}；
\item \dw{社会关系高度和谐，人们精神境界极大提高}；
\item \dw{每个人自由而全面的发展，人类从必然王国向自由王国的飞跃}。
\end{enumerate}

三条的次序也有讲究：先是物质条件（财富与分配），再是社会关系与精神面貌，最后落到人的发展——这正是马克思主义的最高价值目标。

\subhdr{8.5\quad “两个必然”：趋势判断}

《共产党宣言》的结论是：“资产阶级的灭亡和无产阶级的胜利是同样不可避免的。”这就是\dw{两个必然}：资本主义必然灭亡，社会主义必然胜利。它回答的是\dw{方向}问题——历史的总趋势。

\subhdr{8.6\quad “两个绝不会”：条件判断，以及三篇文献的对应}

《〈政治经济学批判〉序言》（1859）写道：

\quoteorig{“无论哪一个社会形态，在它所能容纳的全部生产力发挥出来以前，是决不会灭亡的；而新的更高的生产关系，在它的物质存在条件在旧社会的胎胞里成熟以前，是决不会出现的。”}

这就是\dw{两个绝不会}：旧社会形态不会在生产力还有发展余地时灭亡，新社会形态不会在其物质条件成熟前出现。它回答的是\dw{条件}问题——变革的长期性与条件性。

\textbf{两者为什么不矛盾。}“两个必然”讲的是历史总趋势（生产力与生产关系的矛盾运动最终指向社会主义），“两个绝不会”讲的是这一趋势实现的条件与节奏（生产力没有充分发展，旧制度不会自动退出；新制度的物质条件不成熟，也不会凭空出现）。前者是\dw{必然性}，后者是\dw{条件性}；合起来才是完整的唯物史观：既反对“资本主义永世长存”的幻想，也反对“急于过渡”的空想。

\textbf{小例（鞋与脚）。}把社会形态比作鞋、生产力比作脚：脚长大到鞋装不下时，鞋一定要换（两个必然）；但鞋没磨破、脚还没长到那个程度时，鞋不会自己掉下来（两个绝不会）。这个比喻只用于理解“趋势”与“条件”的关系，不能代替生产力与生产关系矛盾运动的论证。

\textbf{三篇文献的时间线（高频考点）。}

\begin{center}\small
\begin{tabular}{@{}p{3.0cm}p{5.6cm}p{6.4cm}@{}}
\truletop
\thead{文献} & \thead{时间与地位} & \thead{核心论断}\\
\trulemid
《共产党宣言》 & 1848，马克思主义公开问世 & 两个必然；阶级斗争学说；无产阶级的历史使命\\
\trulemid
《〈政治经济学批判〉序言》 & 1859，唯物史观的经典表述 & 两个绝不会；生产力决定生产关系、经济基础决定上层建筑\\
\trulemid
《资本论》第一卷 & 1867，剩余价值学说的系统阐述 & 剩余价值的生产；资本主义积累的一般规律\\
\trulemid
《哥达纲领批判》 & 1875，共产主义形态学说 & 共产主义社会分为第一阶段和高级阶段；第一阶段按劳分配，高级阶段按需分配\\
\trulemid
《帝国主义是资本主义的最高阶段》 & 1916，垄断资本主义理论 & 五个基本经济特征；帝国主义的历史地位\\
\trulebot
\end{tabular}
\end{center}

第 15 题的错法是把“两个绝不会”记成《共产党宣言》。一个简单的分辨办法：\dw{《宣言》讲“必然”（灭亡与胜利不可避免），《序言》讲“绝不会”（不会提前灭亡、不会提前出现）}——一个讲方向，一个讲条件，出处自然分开。

\subhdr{8.7\quad 本节排雷}

\begin{center}\small
\begin{tabular}{@{}p{6.2cm}p{9.0cm}@{}}
\truletop
\thead{错误说法} & \thead{正确表述}\\
\trulemid
“两个绝不会出自《共产党宣言》” & 出自《〈政治经济学批判〉序言》\\
\trulemid
“两个必然而两个绝不会相互矛盾” & 前者讲趋势，后者讲条件，统一于唯物史观\\
\trulemid
“共产主义第一阶段实行按需分配” & 第一阶段（社会主义）实行按劳分配，高级阶段才按需分配\\
\trulemid
“空想社会主义是科学社会主义的一个流派” & 空想社会主义是理论来源；科学社会主义是组成部分\\
\trulemid
“经济文化落后国家可以跳过生产力的发展” & 恰恰相反：必须大力发展生产力，因而建设是长期、艰巨的\\
\trulebot
\end{tabular}
\end{center}

\hdr{§9}{三张速查卡：定位词、固定条目、排雷词}

考前最后一遍只看这三张卡。三张卡分别对应 §1 的四条判据：卡 A 管定位词，卡 B 管条目数，卡 C 管定语替换。

\subhdr{卡 A\quad 定位词与“第一观点”}

\begin{center}\small
\begin{tabular}{@{}p{4.0cm}p{7.2cm}p{4.0cm}@{}}
\truletop
\thead{问什么} & \thead{答什么} & \thead{备注}\\
\trulemid
三大规律各揭示什么 & 对立统一：源泉和动力；质量互变：形式和状态；否定之否定：方向和道路 & 第一章第三节\\
\trulemid
唯物辩证法的根本观点 & 矛盾的观点（因为对立统一是实质和核心） & 第一章\\
\trulemid
认识论首要的基本观点 & 实践的观点 & 第二章\\
\trulemid
辩证法的实质和核心 & 对立统一规律 & 第一章\\
\trulemid
矛盾问题的精髓 & 共性与个性、绝对与相对的关系 & 第一章\\
\trulemid
政治经济学的枢纽 & 劳动二重性学说 & 第四章\\
\trulemid
革命的根本问题 & 国家政权问题 & 第三章（第六章同结论）\\
\trulemid
货币转化为资本的前提 & 劳动力成为商品 & 第四章\\
\trulemid
资本主义工资的本质 & 劳动力价值或价格的转化形式 & 第四章\\
\trulemid
国家垄断资本主义的实质 & 资产阶级国家力量同垄断组织力量的结合，性质未变 & 第五章\\
\trulebot
\end{tabular}
\end{center}

\subhdr{卡 B\quad 固定条目（多选题按条目数核对）}

\begin{center}\small
\begin{tabular}{@{}p{4.2cm}p{11.0cm}@{}}
\truletop
\thead{条目} & \thead{内容}\\
\trulemid
劳动力商品价值三要素 & 本人生活资料；家属生活资料；教育培训费用\\
\trulemid
劳动力成为商品两条件 & 有人身自由；丧失一切生产资料和生活资料\\
\trulemid
列宁五特征 & 垄断组织起决定作用；金融寡头；资本输出；国际垄断同盟；领土瓜分完毕\\
\trulemid
货币五职能 & 价值尺度、流通手段、贮藏手段、支付手段、世界货币\\
\trulemid
实践三特征 & 直接现实性、自觉能动性、社会历史性\\
\trulemid
实践三形式 & 物质生产实践、社会政治实践、科学文化实践\\
\trulemid
社会存在三要素 & 自然地理环境、人口因素、物质生产方式\\
\trulemid
经济全球化四表现 & 生产、贸易、金融、企业经营全球化\\
\trulemid
共产主义三特征 & 财富极大丰富按需分配；社会关系和谐精神境界提高；人的自由全面发展\\
\trulemid
空想社会主义三阶段 & 早期空想；空想平均共产主义；批判的空想社会主义\\
\trulebot
\end{tabular}
\end{center}

\subhdr{卡 C\quad 排雷词（见到就警惕）}

\begin{center}\small
\begin{tabular}{@{}p{5.6cm}p{9.6cm}@{}}
\truletop
\thead{词} & \thead{为什么危险}\\
\trulemid
直观、照镜子、消极被动 & 旧唯物主义认识论的标签\\
\trulemid
先验、生而知之、天赋观念 & 唯心主义先验论\\
\trulemid
改变（经济的资本主义）性质 & 国家垄断资本主义只是局部调整\\
\trulemid
消除／克服（资本主义基本矛盾） & 只能缓和，不能消除\\
\trulemid
一次完成、终极真理 & 违背认识的反复性、无限性；真理的相对性\\
\trulemid
有用即真理 & 实用主义真理观\\
\trulemid
源泉和动力（配否定之否定） & 定位词错配\\
\trulemid
《共产党宣言》＋两个绝不会 & 出处错配，应为《〈政治经济学批判〉序言》\\
\trulebot
\end{tabular}
\end{center}

% ======================= PART8 =======================

\clearpage
\hdr{§10}{自测十题：只测“边界”，不测记忆量}

这十题按你原来的错法设计，材料全部换过，数字与说法都不重复原卷。做题时只允许看题干，不要翻前面的正文；做完再对附录 A 的详解。每题都标注了它考的是哪一条判据。

\practice{P1}{自测}{下列判断中，正确揭示质量互变规律在唯物辩证法体系中地位的是（\quad）\\
A．它揭示了事物发展的源泉和动力\\
B．它揭示了事物发展的形式和状态\\
C．它揭示了事物发展的方向和道路\\
D．它揭示了事物发展的基本趋势和最终结果}{三大规律的定位词（§3.5）}{ans:P1}

\practice{P2}{自测}{“实践的观点是辩证唯物主义认识论首要的和基本的观点。”对这一判断的理解，正确的是（\quad）\\
A．它说明矛盾观点在认识论中不再重要\\
B．它说明认识论与辩证法回答的是同一个问题\\
C．它说明实践是认识的基础，而矛盾观点属于辩证法层次\\
D．它说明认识可以脱离实践而独立发展}{两个“基本观点”的分层（§4.1）}{ans:P2}

\practice{P3}{自测}{某设计团队先让使用者自由通行，再根据踩出的痕迹铺设园区道路。若用马克思主义认识论分析这一过程，下列判断正确的是（\quad）\\
A．这是对事物本来面目的直观反映\\
B．认识是在实践基础上对客体的能动反映\\
C．认识是先于经验的天赋观念展开\\
D．感性认识无需上升为理性认识即可指导实践}{能动反映论与直观反映论的区别（§4.4--4.5）}{ans:P3}

\practice{P4}{自测}{在马克思主义看来，一切革命的根本问题是（\quad）\\
A．统一战线中的领导权问题\\
B．武装斗争的形式问题\\
C．国家政权问题\\
D．依靠群众的程度问题}{革命的根本问题（§5.4）}{ans:P4}

\practice{P5}{自测}{“无论哪一个社会形态，在它所能容纳的全部生产力发挥出来以前，是决不会灭亡的……”这一论断出自（\quad）\\
A．《共产党宣言》\\
B．《〈政治经济学批判〉序言》\\
C．《哥达纲领批判》\\
D．《帝国主义是资本主义的最高阶段》}{经典文献与核心论断的对应（§8.6）}{ans:P5}

\practice{P6}{自测}{关于国家垄断资本主义，下列判断正确的是（\quad）\\
A．它改变了经济的资本主义性质\\
B．它是资本主义生产关系的局部调整，没有改变垄断资本主义的性质\\
C．它已经克服了资本主义基本矛盾\\
D．它使国有经济成为全社会成员共同所有的经济}{国家垄断资本主义的实质（§7.4）}{ans:P6}

\practice{P7}{自测}{劳动力商品的价值包括（\quad）\\
A．维持劳动者本人生存所必需的生活资料的价值\\
B．维持劳动者家属生存所必需的生活资料的价值\\
C．劳动者接受教育和训练所支出的费用\\
D．劳动者在生产中耗费的生产资料的价值}{劳动力商品价值的三要素（§6.4）}{ans:P7}

\practice{P8}{自测}{列宁概括的垄断资本主义基本经济特征包括（\quad）\\
A．垄断组织在经济生活中起决定作用\\
B．资本输出具有特别重要的意义\\
C．在金融资本基础上形成金融寡头的统治\\
D．商品输出取代资本输出成为主要扩张形式}{列宁五特征（§7.3）}{ans:P8}

\practice{P9}{自测}{实践对认识的决定作用表现在（\quad）\\
A．实践是认识的来源\\
B．实践是认识发展的动力\\
C．实践是检验认识真理性的唯一标准\\
D．实践是认识的目的}{实践对认识的决定作用（§4.2）}{ans:P9}

\practice{P10}{自测}{关于马克思主义的理论来源与组成部分，下列判断正确的有（\quad）\\
A．德国古典哲学是马克思主义哲学的直接理论来源\\
B．英国古典政治经济学提供了劳动价值论，成为马克思主义政治经济学的来源\\
C．空想社会主义是马克思主义的三个组成部分之一\\
D．科学社会主义是马克思主义的核心与结论}{理论来源与组成部分的区别（§2.2--2.4）}{ans:P10}

\clearpage
\appendix
\hdr{附录 A}{十题详解}

\answerstart{P1}{自测}{判断哪一项正确揭示质量互变规律的地位。选项 A 源泉和动力；B 形式和状态；C 方向和道路；D 基本趋势和最终结果。}

\textbf{答案：B。}质量互变规律揭示事物发展的\dw{形式和状态}：发展采取“渐进—飞跃”的形式，表现为“量变—质变”的状态。A 是对立统一规律的定位词；C 是否定之否定规律的定位词；D 是自造的说法，教材没有这个提法，属于“看起来像”的干扰项。

\textbf{判据回顾。}见 §3.5 的对照表：三条规律回答“为什么发展／怎样发展／向何处发展”，定位词依次是源泉动力、形式状态、方向道路。做题时先判断题干问的是哪一问，再对号。

\answerstart{P2}{自测}{理解“实践的观点是辩证唯物主义认识论首要的和基本的观点”。选项 A 矛盾观点不再重要；B 认识论与辩证法回答同一问题；C 实践是认识的基础，矛盾观点属于辩证法层次；D 认识可脱离实践独立发展。}

\textbf{答案：C。}这一判断的根据是：实践是认识的来源、动力、目的和检验标准，所以在认识论中，实践观点处于首要的、基本的地位。而矛盾的观点是唯物辩证法的根本观点，它属于第一章的层次，并不因为实践观点“首要”就失去重要性，所以 A 错；认识论与辩证法回答的问题不同（一个问认识怎样产生和发展，一个问世界怎样联系和发展），B 错；D 与“实践是认识的基础”直接冲突。

\textbf{判据回顾。}见 §4.1 的两行对照：两个“根本／首要”分属不同层次，题干问哪一层，就答哪一层。

\answerstart{P3}{自测}{设计团队先让使用者自由通行，再按痕迹铺路。选项 A 直观反映；B 在实践基础上对客体的能动反映；C 天赋观念展开；D 感性认识无需上升为理性认识。}

\textbf{答案：B。}材料里的认识过程是：先有目的（路径要服务于行走），再通过实践获取材料（观察痕迹），经思维加工形成方案（比较、概括、设计），最后回到实践检验（铺路并调整）。这既有对对象的摹写，也有主体的创造，正是“在实践基础上对客体的能动反映”。

\textbf{为什么其余三项错。}A 的“直观反映”是旧唯物主义认识论的标签——把认识看成照镜子式的被动复写，恰好抹掉了材料里最关键的“实践—认识—再实践”；C 是唯心主义先验论；D 割裂了感性认识与理性认识，材料中的“按痕迹概括出方案”本身就是从感性材料上升到理性认识的过程。

\answerstart{P4}{自测}{一切革命的根本问题是什么。选项 A 统一战线领导权；B 武装斗争形式；C 国家政权；D 依靠群众程度。}

\textbf{答案：C。}教材讲社会革命时指出：革命的根本问题是国家政权问题。理由是：革命要改变旧的生产关系，而旧的生产关系靠旧的国家政权维护；不夺取政权，就不能改变生产关系，也就没有社会形态的更替。列宁在《论两个政权》中也说“一切革命的根本问题是国家政权问题”。

\textbf{为什么 A 错。}统一战线的领导权是无产阶级在革命同盟军中的关键问题，属于策略与领导权层面；它重要，但不是“根本问题”这一命题的答案。B、D 同理：武装斗争是形式，依靠群众是路线，都不是“根本问题”。

\answerstart{P5}{自测}{“无论哪一个社会形态，在它所能容纳的全部生产力发挥出来以前，是决不会灭亡的……”出自哪篇文献。}

\textbf{答案：B，《〈政治经济学批判〉序言》。}这段话就是“两个绝不会”，出自 1859 年的《〈政治经济学批判〉序言》，它同时是唯物史观的经典表述（生产力决定生产关系、经济基础决定上层建筑）。

\textbf{区分办法。}《共产党宣言》（1848）给的是“两个必然”（资产阶级的灭亡和无产阶级的胜利是同样不可避免的），讲\dw{方向}；《序言》给的是“两个绝不会”，讲\dw{条件}。《哥达纲领批判》讲共产主义两阶段，《帝国主义是资本主义的最高阶段》讲五特征。

\answerstart{P6}{自测}{关于国家垄断资本主义的正确判断。选项 A 改变了经济的资本主义性质；B 是资本主义生产关系的局部调整，没有改变垄断资本主义的性质；C 已克服基本矛盾；D 国有经济成为全社会成员共同所有。}

\textbf{答案：B。}国家垄断资本主义是国家政权与私人垄断资本融合的垄断资本主义，它\dw{没有改变垄断资本主义的性质}，只是资本主义生产关系的局部调整。三条论证：所有制没变（国有企业属于资产阶级国家）、目的没变（服务垄断资本的整体利益）、结果没变（剩余价值规律与基本矛盾仍在起作用）。

A 把“局部调整”读成“性质改变”；C 把“缓和”读成“克服”；D 把“国有”读成“公有”——国有企业的所有权属于资产阶级国家，不等于全社会成员共同所有。

\answerstart{P7}{自测}{劳动力商品的价值包括哪些。选项 A 本人生活资料；B 家属生活资料；C 教育培训费用；D 生产中耗费的生产资料的价值。}

\textbf{答案：A、B、C。}劳动力商品的价值由生产、发展、维持和延续劳动力所必需的生活必需品的价值决定，包括三项：本人生活资料的价值、家属生活资料的价值、教育培训费用；此外还包含历史和道德的因素。

D 不属于：生产资料的价值是\dw{不变资本}的组成部分，它在生产中只被转移、不被创造，与劳动力商品的价值无关。这一项正是“把 c 混进 v”的典型干扰。

\textbf{记忆锚。}活着—繁衍—会干。三项缺一不可：少第一项活不下去，少第二项劳动力无法代际接续，少第三项无法适应需要技能的生产。

\answerstart{P8}{自测}{列宁概括的垄断资本主义基本经济特征包括哪些。选项 A 垄断组织起决定作用；B 资本输出具有特别重要的意义；C 金融寡头的统治；D 商品输出取代资本输出成为主要扩张形式。}

\textbf{答案：A、B、C。}列宁的五特征是：（1）生产和资本的集中造成在经济生活中起决定作用的垄断组织；（2）银行资本与工业资本融合，在金融资本基础上形成金融寡头；（3）资本输出有了特别重要的意义；（4）瓜分世界的资本家国际垄断同盟形成；（5）最大资本主义大国已把世界领土瓜分完毕。

D 把方向说反了：自由竞争阶段以\dw{商品输出}为主，垄断阶段\dw{资本输出}具有特别重要的意义。这一项正是利用“商品输出／资本输出”的易混点设的陷阱。

\answerstart{P9}{自测}{实践对认识的决定作用表现在哪些方面。选项 A 来源；B 动力；C 检验标准；D 目的。}

\textbf{答案：A、B、C、D 全选。}实践是认识的来源、发展的动力、检验真理性的唯一标准、认识的目的——四条构成“实践是认识的基础”的完整内容。四项都是教材原文的表述，没有偷换。

\textbf{提示。}这类“四条全对”的多选题容易让人怀疑“是不是太简单”，但教材的固定条目就是四条，全选是正确操作。判断依据不是“选项看起来都挺好”，而是“这四条确实都在教材的同一段里”。

\answerstart{P10}{自测}{关于理论来源与组成部分的判断。选项 A 德国古典哲学是马克思主义哲学的直接理论来源；B 英国古典政治经济学提供了劳动价值论；C 空想社会主义是三个组成部分之一；D 科学社会主义是核心与结论。}

\textbf{答案：A、B、D。}A、B 说的是“来源”，表述正确；D 说的是“组成部分”的地位，也正确。C 错在把\dw{来源}当成了\dw{组成部分}：空想社会主义是理论来源之一（英法空想社会主义），组成部分是科学社会主义。

\textbf{一句话分辨。}问“从哪儿来”答三来源（德国古典哲学、英国古典政治经济学、英法空想社会主义）；问“由什么构成”答三组成（马克思主义哲学、马克思主义政治经济学、科学社会主义）。

\clearpage
\hdr{附录 B}{依据、范围与未解决项}

\textbf{教材依据。}本讲义按高等教育出版社《马克思主义基本原理》（马克思主义理论研究和建设工程重点教材，2023 年版）的章节体系组织：导论，第一章世界的物质性及发展规律，第二章实践与认识及其发展规律，第三章人类社会及其发展规律，第四章资本主义的本质及规律，第五章资本主义的发展及其趋势，第六章社会主义的发展及其规律，第七章共产主义崇高理想及其最终实现。

\textbf{关于页码。}本次没有获得该教材的电子版或纸质扫描件，因此正文只标章节位置，不标页码；引号内的经典原文（列宁五特征、“两个绝不会”、马克思关于建筑师与蜜蜂的论述）按通行版本核对，若你手上的版本文字略有出入，以教材为准，结论不受影响。

\textbf{覆盖范围。}你列出的九处失分点全部有独立小节（§2.2--2.4、§3.2--3.7、§4.1--4.5、§5.4、§6.2--6.4、§7.3--7.4、§8.6），并补了三处相邻考点：剩余价值生产方法（§6.5）、经济全球化表现（§7.5）、货币转化为资本的前提（§6.3）。为让这些点站得住，另补了必要的体系骨架：三大规律内容（§3）、实践与真理（§4）、唯物史观两大规律（§5.2）、资本积累与周转（§6.6--6.7）、社会再生产（§6.8）、当代资本主义新变化（§7.6）、共产主义特征（§8.4）。

\textbf{未解决项（如实说明）。}

\begin{enumerate}[label=\arabic*.,leftmargin=2.0em]
\item 你的原卷题干与选项文本我没有拿到，只有你转述的错项内容（如“原选 D”“误选 A”）。所以正文按\dw{错法类型}讲评，没有逐字复现原题；如果你把原卷拍照给我，可以再补一份逐题讲评。
\item 教材版次按 2023 年版处理。若你的课程使用 2021 年版或校本讲义，个别表述（如第七章的编排、社会主义一般原则的条数）可能不同，但本讲义涉及的九处考点在两版中结论一致。
\item 讲义已覆盖不等于已掌握。按你当前的学习模式（讲义驱动、考前集中闭卷验收），这些内容要等换题换数字的闭卷验收才能记为“能独立完成”。
\end{enumerate}

\textbf{下一步建议。}先把 §9 的三张卡默写一遍（尤其卡 B 的十条固定条目），再做附录 A 的十题；错的那一题回到对应小节，只读那一节，不要重读全篇。

\end{document}
